\documentclass[a4paper,14pt]{extarticle}

\usepackage{fontspec}
\usepackage{polyglossia}
\setmainlanguage{russian}
\setotherlanguage{english}
\setmainfont[Ligatures=TeX]{Times New Roman}
\newfontfamily\cyrillicfont[Ligatures=TeX]{Times New Roman}
\newfontfamily\cyrillicfonttt{Menlo}[Scale=0.88]
\newfontfamily\cyrillicfontsf{Arial}

\usepackage[
  left=30mm,
  right=15mm,
  top=20mm,
  bottom=20mm
]{geometry}
\usepackage{setspace}
\onehalfspacing
\setlength{\parindent}{1.25cm}
\setlength{\parskip}{0pt}

% Сдаваемая версия ВКР не должна разрывать слова переносами вида
% "обозначе- / ния", особенно на стыке страниц.
\hyphenpenalty=10000
\exhyphenpenalty=10000
\brokenpenalty=10000
\doublehyphendemerits=1000000
\finalhyphendemerits=1000000

\usepackage{indentfirst}
\usepackage{microtype}
\usepackage{graphicx}
\usepackage{adjustbox}
\usepackage{float}
\usepackage{flafter}
\usepackage{caption}
\usepackage{longtable}
\usepackage{booktabs}
\usepackage{array}
\usepackage{calc}
\usepackage{ragged2e}
\usepackage{etoolbox}
\usepackage{enumitem}
\usepackage{amsmath}
\usepackage{amssymb}
\usepackage{fvextra}
\usepackage{listings}
\usepackage{lastpage}
\usepackage{titlesec}
\usepackage{tocloft}
\usepackage{needspace}
\usepackage{placeins}
\usepackage{hyperref}
\usepackage{bookmark}
\usepackage{xcolor}

\graphicspath{{./}{../}{docx_pipeline/figures/png/}{docx_pipeline/figures/svg/}}

\hypersetup{
  unicode=true,
  colorlinks=false,
  pdfborder={0 0 0},
  pdftitle={Методика адаптивного автообзвона с использованием больших языковых моделей и анализа эмоционального состояния собеседника},
  pdfauthor={Ларионов Сергей Иванович}
}

\setcounter{secnumdepth}{0}
\setcounter{tocdepth}{3}

\titleformat{\section}
  {\normalfont\bfseries}
  {}
  {0pt}
  {}
\titlespacing*{\section}{1.25cm}{0pt}{12pt}

\titleformat{\subsection}
  {\normalfont\bfseries}
  {}
  {0pt}
  {}
\titlespacing*{\subsection}{1.25cm}{12pt}{6pt}

\titleformat{\subsubsection}
  {\normalfont\bfseries}
  {}
  {0pt}
  {}
\titlespacing*{\subsubsection}{1.25cm}{10pt}{6pt}

\pretocmd{\section}{\FloatBarrier\clearpage}{}{}
\pretocmd{\subsection}{\Needspace{6\baselineskip}}{}{}
\pretocmd{\subsubsection}{\Needspace{5\baselineskip}}{}{}

\newcommand{\VkrStructuralSection}[1]{%
  \FloatBarrier
  \clearpage
  \phantomsection
  \addcontentsline{toc}{section}{#1}%
  \begin{center}\bfseries #1\end{center}%
  \par\nopagebreak\vspace{12pt}%
}

\newcommand{\VkrStructuralSectionNoToc}[1]{%
  \begin{center}\bfseries #1\end{center}%
  \par\nopagebreak\vspace{12pt}%
}

\newcommand{\VkrAppendixSection}[2]{%
  \FloatBarrier
  \clearpage
  \phantomsection
  \addcontentsline{toc}{section}{ПРИЛОЖЕНИЕ #1 #2}%
  \begin{center}\bfseries ПРИЛОЖЕНИЕ #1\end{center}%
  \par\nopagebreak\vspace{6pt}%
  \begin{center}\bfseries #2\end{center}%
  \par\nopagebreak\vspace{12pt}%
}

\renewcommand{\cftsecleader}{\cftdotfill{\cftdotsep}}
\makeatletter
\renewcommand{\tableofcontents}{%
  \section*{СОДЕРЖАНИЕ}%
  \@starttoc{toc}%
}
\makeatother
\setlength{\cftbeforesecskip}{2pt}
\setlength{\cftbeforesubsecskip}{1pt}

\captionsetup{
  font={normalsize,stretch=1},
  labelformat=empty,
  labelsep=none,
  justification=centering,
  singlelinecheck=false,
  skip=6pt
}
\captionsetup[lstlisting]{
  font={normalsize,stretch=1},
  labelformat=simple,
  labelsep=endash,
  justification=raggedright,
  singlelinecheck=false,
  skip=6pt
}

\setlist[itemize]{label=\textemdash,leftmargin=1.25cm,itemsep=0pt,topsep=0pt,parsep=0pt}
\setlist[enumerate]{leftmargin=1.25cm,itemsep=0pt,topsep=0pt,parsep=0pt}

\lstset{
  basicstyle=\ttfamily\small,
  breaklines=true,
  columns=fullflexible,
  keepspaces=true,
  frame=single,
  framerule=0.4pt,
  framesep=4pt,
  captionpos=t,
  rulecolor=\color{black},
  xleftmargin=0pt,
  xrightmargin=0pt,
  aboveskip=6pt,
  belowskip=6pt
}
\renewcommand{\lstlistingname}{Листинг}
\RecustomVerbatimEnvironment{verbatim}{Verbatim}{
  breaklines=true,
  breakanywhere=true,
  fontsize=\small,
  frame=single,
  framerule=0.4pt,
  framesep=4pt
}
\renewcommand{\texttt}[1]{\textnormal{#1}}

\sloppy
\emergencystretch=2em
\clubpenalty=10000
\widowpenalty=10000
\displaywidowpenalty=10000
\pagestyle{plain}

\providecommand{\tightlist}{%
  \setlength{\itemsep}{0pt}\setlength{\parskip}{0pt}}
\newcommand{\VkrFigureFull}[1]{%
  \adjustbox{max width=\linewidth,max totalheight=0.74\textheight,keepaspectratio,center}{#1}%
}
\newcommand{\VkrFigureInline}[1]{%
  \adjustbox{max width=0.94\linewidth,max totalheight=0.55\textheight,keepaspectratio,center}{#1}%
}
\newcommand{\VkrFigureArchitecture}[1]{%
  \adjustbox{max width=0.98\linewidth,max totalheight=0.74\textheight,keepaspectratio,center}{#1}%
}
\newcommand{\VkrFigureTurnCycle}[1]{%
  \adjustbox{max width=0.98\linewidth,max totalheight=0.76\textheight,keepaspectratio,center}{#1}%
}
\newcommand{\VkrFigureStagePolicy}[1]{%
  \adjustbox{max width=0.98\linewidth,max totalheight=0.65\textheight,keepaspectratio,center}{#1}%
}
\newcommand{\VkrFigureExperimentDesign}[1]{%
  \adjustbox{max width=0.98\linewidth,max totalheight=0.60\textheight,keepaspectratio,center}{#1}%
}
\newcommand{\VkrFigureTall}[1]{%
  \adjustbox{max width=0.88\linewidth,max totalheight=0.50\textheight,keepaspectratio,center}{#1}%
}
\newcommand{\VkrFigureVertical}[1]{%
  \adjustbox{max width=0.88\linewidth,max totalheight=0.56\textheight,keepaspectratio,center}{#1}%
}
\newcommand{\VkrFigureCompact}[1]{%
  \adjustbox{max width=0.92\linewidth,max totalheight=0.48\textheight,keepaspectratio,center}{#1}%
}
\newcommand{\VkrFigureWideCompact}[1]{%
  \adjustbox{max width=0.88\linewidth,max totalheight=0.38\textheight,keepaspectratio,center}{#1}%
}
\providecommand{\pandocbounded}[1]{\VkrFigureFull{#1}}
\providecommand{\passthrough}[1]{#1}
\newcounter{none}

\setlength{\LTpre}{0pt}
\setlength{\LTpost}{6pt}
\setlength{\tabcolsep}{3.5pt}
\setlength{\extrarowheight}{2pt}
\renewcommand{\arraystretch}{1.22}
\AtBeginEnvironment{longtable}{\small\singlespacing}

\newcommand{\VkrTableCaption}[1]{%
  \Needspace{0.44\textheight}%
  \par\smallskip
  {\singlespacing\noindent #1\par}%
  \vspace{2pt}%
}


\begin{document}
\hypersetup{pageanchor=false}
\newcommand{\VkrTitleHeaderLine}[1]{%
  \noindent{\fontsize{12pt}{14pt}\selectfont\makebox[\textwidth][c]{#1}\par}%
}
\newcommand{\VkrUdkField}{УДК \underline{\hspace{1.8cm}}}

\begin{titlepage}
\fontsize{12pt}{14pt}\selectfont
\singlespacing
\VkrTitleHeaderLine{МИНИСТЕРСТВО НАУКИ И ВЫСШЕГО ОБРАЗОВАНИЯ РОССИЙСКОЙ ФЕДЕРАЦИИ}

\vspace{0.95cm}

\VkrTitleHeaderLine{НАЦИОНАЛЬНЫЙ ИССЛЕДОВАТЕЛЬСКИЙ ЯДЕРНЫЙ УНИВЕРСИТЕТ «МИФИ»}

\vspace{0.95cm}

\VkrTitleHeaderLine{ИНСТИТУТ ИНТЕЛЛЕКТУАЛЬНЫХ КИБЕРНЕТИЧЕСКИХ СИСТЕМ}

\vspace{0.95cm}

\VkrTitleHeaderLine{КАФЕДРА № 42 «КРИПТОЛОГИЯ И КИБЕРБЕЗОПАСНОСТЬ»}

\vspace{0.9cm}

\begin{flushright}
На правах рукописи

\vspace{0.2cm}

\VkrUdkField
\end{flushright}

\vspace{0.9cm}

\begin{center}
Ларионов Сергей Иванович

\vspace{0.85cm}

\textbf{«Методика адаптивного автообзвона с использованием больших языковых моделей и анализа эмоционального состояния собеседника»}

\vspace{1.0cm}

Выпускная квалификационная работа магистра

\vspace{0.75cm}

Направление подготовки 01.04.02\\
«Машинное обучение»
\end{center}

\vspace{0.55cm}

\begin{flushright}
\begin{minipage}{0.46\textwidth}
Выпускная квалификационная\\
работа защищена

\vspace{0.15cm}

«\underline{\hspace{0.8cm}}»\underline{\hspace{2.5cm}}2026 г.

\vspace{0.18cm}

Оценка

\vspace{0.12cm}

\underline{\hspace{3.0cm}}

\vspace{0.18cm}

Секретарь ГЭК \underline{\hspace{2.2cm}}
\end{minipage}
\end{flushright}

\vfill

\begin{center}
г. Москва\\
2026
\end{center}
\end{titlepage}

\begin{titlepage}
\fontsize{12pt}{14pt}\selectfont
\singlespacing
\VkrTitleHeaderLine{МИНИСТЕРСТВО НАУКИ И ВЫСШЕГО ОБРАЗОВАНИЯ РОССИЙСКОЙ ФЕДЕРАЦИИ}

\vspace{0.95cm}

\VkrTitleHeaderLine{НАЦИОНАЛЬНЫЙ ИССЛЕДОВАТЕЛЬСКИЙ ЯДЕРНЫЙ УНИВЕРСИТЕТ «МИФИ»}

\vspace{0.95cm}

\VkrTitleHeaderLine{ИНСТИТУТ ИНТЕЛЛЕКТУАЛЬНЫХ КИБЕРНЕТИЧЕСКИХ СИСТЕМ}

\vspace{0.95cm}

\VkrTitleHeaderLine{КАФЕДРА № 42 «КРИПТОЛОГИЯ И КИБЕРБЕЗОПАСНОСТЬ»}

\vspace{0.9cm}

\begin{flushright}
На правах рукописи

\vspace{0.2cm}

\VkrUdkField
\end{flushright}

\vspace{0.9cm}

\begin{center}
Ларионов Сергей Иванович

\vspace{0.85cm}

\textbf{«Методика адаптивного автообзвона с использованием больших языковых моделей и анализа эмоционального состояния собеседника»}

\vspace{1.0cm}

Выпускная квалификационная работа магистра

\vspace{0.75cm}

Направление подготовки 01.04.02\\
«Машинное обучение»
\end{center}

\vfill

\begin{center}
г. Москва\\
2026
\end{center}
\end{titlepage}

\setcounter{page}{2}
\hypersetup{pageanchor=true}

\VkrStructuralSectionNoToc{РЕФЕРАТ}\label{ux440ux435ux444ux435ux440ux430ux442}

Работа содержит \pageref{LastPage} с., 7 рис., 13 табл., 22 источн., 1
прил.

АДАПТИВНЫЙ АВТООБЗВОН, БОЛЬШИЕ ЯЗЫКОВЫЕ МОДЕЛИ, РАСПОЗНАВАНИЕ ЭМОЦИЙ В
РЕЧИ, ГОЛОСОВОЙ АГЕНТ, СЛОЙ ПРАВИЛ АДАПТАЦИИ, ДИАЛОГОВАЯ СИСТЕМА,
A/B-ЭКСПЕРИМЕНТ, КОНТАКТ-ЦЕНТР.

Объектом исследования являются процессы автоматизированного исходящего
телефонного взаимодействия с клиентами. Цель работы состоит в разработке
воспроизводимой методики адаптивного автообзвона, обеспечивающей
управляемую корректировку сценария диалога на основе распознавания
эмоционального состояния собеседника и генерации реплик с использованием
больших языковых моделей.

В работе использованы анализ литературы по emotion-aware dialogue
systems, speech emotion recognition и voice-пайплайнам,
структурно-функциональное проектирование диалоговой системы, ручная
экспертная оценка транскриптов звонков по фиксированной схеме
кодирования экспертных оценок и контролируемое A/B-сравнение
\texttt{baseline}- и \texttt{adaptive}-сценариев.

Результатом работы является формализованная методика адаптивного
автообзвона: архитектура конвейера
\texttt{VAD}--\texttt{ASR}--\texttt{SER}--normalization layer--policy
layer--\texttt{LLM}--\texttt{TTS}, правила интерпретации эмоций, учет
\texttt{confidence}, требования к журналированию и протокол
экспериментальной проверки. При доле недозвонов около \(65\%\) для
анализа была сформирована выборка звонков длительностью более \(40\)
секунд; она включила 267 экспериментальных звонков. В primary-срезе
adaptive-группа показала более высокую долю целевых исходов:
\(25{,}2\%\) против \(17{,}9\%\) у baseline-группы, прирост составил
\(7{,}36\) процентного пункта. Результат трактуется как положительный
пилотный сигнал без статистически доказанного превосходства на уровне
\(0{,}05\); практическая область применения связана с построением
русскоязычных голосовых агентов для контакт-центров и сервисных
обзвонов.

\clearpage

\VkrStructuralSectionNoToc{СОДЕРЖАНИЕ}
\makeatletter\@starttoc{toc}\makeatother
\clearpage

\VkrStructuralSection{ТЕРМИНЫ И ОПРЕДЕЛЕНИЯ}\label{ux442ux435ux440ux43cux438ux43dux44b-ux438-ux43eux43fux440ux435ux434ux435ux43bux435ux43dux438ux44f}

В настоящем отчете применяют следующие термины с соответствующими
определениями.

Адаптивный автообзвон --- автоматизированный исходящий телефонный
диалог, в котором система изменяет тактику ответа в зависимости от
текущего состояния собеседника, стадии сценария и внутренних правил
принятия решения.

Методика адаптивного автообзвона --- воспроизводимое описание состава
модулей, входов, выходов, правил адаптации, ограничений и процедур
логирования, обеспечивающее реализацию и экспериментальную проверку
адаптивного автообзвона.

Сценарий диалога --- регламентированная последовательность стадий и
допустимых переходов исходящего разговора, определяющая цель каждого
шага, ограничения на формулировки и набор допустимых исходов.

Диалоговый такт --- минимальный цикл обработки в разговорной системе,
включающий прием очередной пользовательской реплики, распознавание речи
и эмоции, выбор policy-действия, генерацию ответа и запись событий в
журнал.

Эмоциональное состояние собеседника --- нормализованная интерпретация
текущего эмоционального режима абонента, используемая системой для
выбора следующей тактики разговора.

Нормализация эмоций --- процедура отображения первичных выходов модели
\texttt{SER} в ограниченный рабочий набор состояний методики, например в
нейтральность, заинтересованность, раздражение, злость и
неопределенность.

Confidence эмоциональной гипотезы --- числовая оценка надежности
распознанного эмоционального состояния, используемая как управляющий
параметр при выборе допустимой глубины адаптации.

Policy layer --- детерминированный слой принятия решений, который
преобразует состояние собеседника, значение \texttt{confidence}, стадию
сценария и историю диалога в следующее действие системы.

Prompt builder --- модуль, преобразующий решение \texttt{policy\ layer}
в структурированную инструкцию для языковой модели с учетом сценарной
цели, ограничений безопасности и контекста разговора.

Базовый сценарий (\texttt{baseline}) --- контрольный вариант системы, в
котором используется тот же сценарный каркас и тот же технологический
контур, но без изменения тактики по данным \texttt{SER}.

Адаптивный сценарий (\texttt{adaptive}) --- экспериментальный вариант
системы, в котором тактика следующего ответа изменяется на основе
эмоционального состояния собеседника, значения \texttt{confidence} и
стадии разговора.

Fallback-режим --- безопасный режим работы системы, применяемый при
низкой надежности эмоциональной гипотезы, нарушении ограничений ответа,
технической неопределенности либо иных условиях, при которых агрессивная
адаптация недопустима.

Эскалация на оператора --- управляемый перевод разговора от
автоматической системы к человеку-оператору при наступлении заранее
определенных условий, например при высокой вероятности конфликта.

Целевое действие --- заранее операционализированный исход разговора,
свидетельствующий о достижении основной цели сценария, например о
согласовании следующего содержательного контакта.

Квалифицированный лид --- контакт, по которому в ходе разговора
подтверждена минимальная релевантность предложения и возможность
дальнейшего взаимодействия, даже если целевое действие еще не
зафиксировано.

Журналирование такта --- обязательная запись входных сигналов, выходов
модулей, policy-решения, сгенерированной реплики и исхода шага для
последующего анализа и воспроизводимости.

Воспроизводимость методики --- свойство исследования, при котором другой
разработчик или исследователь может восстановить архитектуру, правила
адаптации, экспериментальный контур и критерии оценки без обращения к
неформализованным эвристикам.

Управляемость системы --- способность ограничивать поведение языковой
модели рамками сценария, явными policy-правилами и процедурами
валидации, исключающими неконтролируемые критические переходы.

Контрольная группа \texttt{A} --- совокупность звонков, обрабатываемых в
режиме \texttt{baseline}, то есть без эмоциональной адаптации на уровне
\texttt{policy\ layer}.

Экспериментальная группа \texttt{B} --- совокупность звонков,
обрабатываемых в режиме \texttt{adaptive}, то есть с использованием
\texttt{SER}, \texttt{confidence} и правил адаптации.

Единица анализа --- один завершенный телефонный контакт, по которому
зафиксирован итог разговора и доступен достаточный лог для анализа.

Конверсионная метрика --- показатель, отражающий достижение целевого
действия или близкого к нему исхода в рамках зафиксированного сценария.

Диалоговая метрика --- показатель, описывающий структуру разговора:
длительность, число реплик, частоту досрочного завершения и иные
свойства взаимодействия.

Эмоционально-поведенческая метрика --- показатель, фиксирующий
негативные реакции, эскалации, конфликтные завершения и иные эффекты,
связанные с состоянием собеседника и тактикой системы.

\VkrStructuralSection{ПЕРЕЧЕНЬ СОКРАЩЕНИЙ И ОБОЗНАЧЕНИЙ}\label{ux43fux435ux440ux435ux447ux435ux43dux44c-ux441ux43eux43aux440ux430ux449ux435ux43dux438ux439-ux438-ux43eux431ux43eux437ux43dux430ux447ux435ux43dux438ux439}

В настоящем отчете применяют следующие сокращения и обозначения.
Основные сокращения приведены в таблице 1, а используемые далее
обозначения методики и эксперимента --- в таблице 2.

\VkrTableCaption{Таблица 1 --- Основные сокращения, используемые в работе}

{\def\LTcaptype{none} % do not increment counter
\begin{longtable}[]{|%
>{\RaggedRight\arraybackslash}p{(\linewidth - 11\tabcolsep) * \real{0.3333}}|
  >{\RaggedRight\arraybackslash}p{(\linewidth - 11\tabcolsep) * \real{0.3333}}|
  >{\RaggedRight\arraybackslash}p{(\linewidth - 11\tabcolsep) * \real{0.3333}}|}
\hline
\begin{minipage}[b]{\linewidth}\RaggedRight
Сокращение
\end{minipage} & \begin{minipage}[b]{\linewidth}\RaggedRight
Расшифровка
\end{minipage} & \begin{minipage}[b]{\linewidth}\RaggedRight
Использование в работе
\end{minipage} \\
\hline
\endfirsthead
\hline\multicolumn{3}{|l|}{\small Продолжение таблицы 1}\\
\hline
\begin{minipage}[b]{\linewidth}\RaggedRight
Сокращение
\end{minipage} & \begin{minipage}[b]{\linewidth}\RaggedRight
Расшифровка
\end{minipage} & \begin{minipage}[b]{\linewidth}\RaggedRight
Использование в работе
\end{minipage} \\
\hline
\endhead
\hline
\endfoot
\hline
\endlastfoot
\texttt{A/B} & Comparative A/B testing & Сравнение контрольного и
экспериментального режимов. \\
\texttt{API} & Application Programming Interface & Интерфейс
взаимодействия программных модулей и внешних сервисов. \\
\texttt{ASR} & Automatic Speech Recognition & Автоматическое
распознавание речи. \\
\texttt{LLM} & Large Language Model & Большая языковая модель,
генерирующая следующую реплику. \\
\texttt{RNNT} & Recurrent Neural Network Transducer & Архитектурный
класс end-to-end моделей распознавания речи. \\
\texttt{SER} & Speech Emotion Recognition & Автоматическое распознавание
эмоций по речи. \\
\texttt{SSL} & Self-Supervised Learning & Подход к обучению речевых
представлений без ручной разметки каждого примера. \\
\texttt{TTS} & Text-to-Speech & Синтез речи по тексту. \\
\texttt{VAD} & Voice Activity Detection & Детекция речевой активности и
сегментация аудиосигнала. \\
\end{longtable}
}

\VkrTableCaption{Таблица 2 --- Основные обозначения, используемые в методике адаптивного автообзвона}

{\def\LTcaptype{none} % do not increment counter
\begin{longtable}[]{|%
>{\RaggedRight\arraybackslash}p{(\linewidth - 9\tabcolsep) * \real{0.5000}}|
  >{\RaggedRight\arraybackslash}p{(\linewidth - 9\tabcolsep) * \real{0.5000}}|}
\hline
\begin{minipage}[b]{\linewidth}\RaggedRight
Обозначение
\end{minipage} & \begin{minipage}[b]{\linewidth}\RaggedRight
Значение
\end{minipage} \\
\hline
\endfirsthead
\hline\multicolumn{2}{|l|}{\small Продолжение таблицы 2}\\
\hline
\begin{minipage}[b]{\linewidth}\RaggedRight
Обозначение
\end{minipage} & \begin{minipage}[b]{\linewidth}\RaggedRight
Значение
\end{minipage} \\
\hline
\endhead
\hline
\endfoot
\hline
\endlastfoot
\(\theta_{\mathrm{low}}\) & Нижний порог надежности эмоциональной
гипотезы. При значениях ниже этого порога система переходит к состоянию
неопределенности. \\
\(\theta_{\mathrm{high}}\) & Верхний порог надежности эмоциональной
гипотезы. Начиная с этого значения допускается полный набор
policy-действий. \\
\(A\) & Контрольная группа эксперимента, соответствующая
\texttt{baseline}-сценарию. \\
\(B\) & Экспериментальная группа эксперимента, соответствующая
\texttt{adaptive}-сценарию. \\
\(c_t\) & Числовая оценка надежности предсказанного эмоционального
состояния на такте диалога. \\
\texttt{baseline} & Вариант системы без эмоциональной адаптации на
уровне \texttt{policy\ layer}. \\
\texttt{adaptive} & Вариант системы с эмоциональной адаптацией на основе
\texttt{SER}, \texttt{confidence} и стадии диалога. \\
\end{longtable}
}

При первом появлении сокращения в основном тексте ВКР целесообразно
приводить его полную расшифровку, а далее использовать краткую форму.
Обозначения \(\theta_{\mathrm{low}}\), \(\theta_{\mathrm{high}}\) и
\(c_t\) в главе 2 интерпретируются не только как формальные символы, но
и как содержательные параметры предложенной методики.

\VkrStructuralSection{ВВЕДЕНИЕ}\label{ux432ux432ux435ux434ux435ux43dux438ux435}

Современные разговорные системы на основе больших языковых моделей и
речевых технологий все активнее внедряются в прикладные контуры
взаимодействия с клиентами. Одним из наиболее практически значимых
направлений такого внедрения является автоматизированный исходящий
обзвон, применяемый в маркетинговых, сервисных и информационных
сценариях. Вместе с тем эффективность автообзвона определяется не только
качеством распознавания речи и корректностью сгенерированного текста, но
и способностью системы учитывать текущее состояние собеседника. В
реальном телефонном диалоге признаки раздражения, заинтересованности,
нейтральности или нарастающего конфликта должны влиять на выбор
дальнейшей тактики разговора. Отсутствие такой адаптации ведет либо к
избыточно жесткому и шаблонному взаимодействию, либо к повышению
вероятности досрочного завершения звонка и конфликтного исхода.

Актуальность темы обусловлена несколькими факторами. Во-первых,
современные emotion-aware исследования показывают, что эмоциональный
контекст следует учитывать не только при стилизации ответа, но и на
уровне выбора следующего действия системы.

Во-вторых, в русскоязычном прикладном контуре сохраняется дефицит
воспроизводимых методик, которые объединяли бы распознавание эмоций в
речи, rule-based policy-адаптацию, генерацию следующей реплики и
требования к экспериментальной проверке.

В-третьих, для задач исходящего обзвона особенно важны объяснимость,
управляемость и безопасность поведения системы, поскольку диалог
разворачивается в регламентированном сценарии и должен соответствовать
ограничениям домена {[}1; 2; 3{]}.

Степень разработанности темы можно охарактеризовать как существенную, но
фрагментированную. В литературе широко представлены исследования по
emotion-aware dialogue, empathetic generation, speech emotion
recognition, spoken dialogue systems и full-duplex voice interaction.
Однако существующие работы либо ориентированы преимущественно на
англоязычный материал, либо рассматривают эмоцию как стилистический
фактор, либо не доводят архитектуру до уровня воспроизводимого
прикладного контура исходящего автообзвона. Тем самым сохраняется
исследовательская ниша для формализации методики, ориентированной на
русский язык, регламентированный телефонный сценарий и controlled
\texttt{baseline\ vs\ adaptive} эксперимент {[}3; 4; 5; 6{]}.

По результатам анализа открытых источников установлено, что в литературе
и доступных прикладных описаниях существуют отдельные решения по
распознаванию эмоций в речи, emotion-aware dialogue и голосовым ботам
для контакт-центров, однако не выявлена завершенная воспроизводимая
методика русскоязычного адаптивного автообзвона, в которой эмоциональное
состояние и надежность его распознавания используются как формальные
управляющие параметры \texttt{policy\ layer}, а эффективность
проверяется в A/B-эксперименте по схеме \texttt{baseline\ vs\ adaptive}.
Данное утверждение относится к корпусу открыто опубликованных материалов
и не исключает существования закрытых корпоративных реализаций,
недоступных для академического анализа; оно сформулировано по итогам
сопоставления работ по emotion-aware dialogue, spoken dialogue и
call-center pipeline {[}1; 2; 4; 7; 8{]}.

Объектом исследования являются процессы автоматизированного исходящего
телефонного взаимодействия с клиентами в задачах маркетингового и
сервисного обзвона. Предмет исследования составляют методы и алгоритмы
динамической адаптации сценария диалога в реальном времени на основе
распознавания эмоционального состояния собеседника и генерации ответов с
использованием больших языковых моделей.

Цель работы состоит в разработке и верификации воспроизводимой методики
адаптивного автообзвона, обеспечивающей управляемую корректировку
сценария разговора по сигналам эмоционального состояния собеседника,
распознаваемым из речи в реальном времени.

Для достижения поставленной цели в работе решаются следующие задачи:

\begin{enumerate}
\def\labelenumi{\arabic{enumi})}
\tightlist
\item
  провести аналитический обзор современных подходов к emotion-aware
  dialogue, speech emotion recognition и архитектурам
  \texttt{ASR\ +\ LLM\ +\ TTS};
\item
  сформировать формализованную методику адаптации диалога для задач
  автообзвона с явными правилами принятия решений;
\item
  обосновать выбор модулей конвейера \texttt{VAD}, \texttt{ASR},
  \texttt{SER}, \texttt{LLM} и \texttt{TTS} для русскоязычного контекста
  и описать требования к их интеграции;
\item
  разработать прототип, реализующий предложенную методику в условиях
  пилотного исходящего обзвона;
\item
  спроектировать и провести A/B-эксперимент
  \texttt{baseline\ vs\ adaptive} с фиксированными метриками
  эффективности;
\item
  оценить ограничения, риски и условия применимости предложенного
  подхода.
\end{enumerate}

Научная новизна работы в рабочей постановке состоит в следующем.
Во-первых, предлагается методика адаптивного автообзвона, в которой
эмоциональное состояние абонента используется не только для изменения
стиля ответа, но и как управляемый сигнал выбора следующей тактики
диалога.

Во-вторых, вводится правило адаптации, связывающее класс эмоции,
значение \texttt{confidence} и стадию диалога с конкретным действием
\texttt{policy\ layer}.

В-третьих, формируется воспроизводимый экспериментальный контур для
количественной проверки вклада эмоциональной адаптации в бизнес- и
диалоговые метрики исходящего обзвона.

В-четвертых, новизна работы состоит не в изолированном использовании
\texttt{SER}, \texttt{LLM} или голосового робота как таковых, а в их
объединении в формальную русскоязычную методику, пригодную для
объяснимой реализации и проверяемую в контролируемом сравнении
\texttt{baseline\ vs\ adaptive}.

Практическая значимость работы состоит в том, что предлагаемая методика
ориентирована на использование в прикладном контуре агентства или
контактного центра и позволяет повысить управляемость поведения
автообзвона, снизить долю конфликтных взаимодействий и увеличить долю
конструктивно завершенных разговоров по сравнению с неадаптивным
сценарием. Отдельная ценность подхода заключается в том, что он
описывается как воспроизводимая инженерная схема, а не как набор
разрозненных эвристик.

Методологическую основу исследования составляют труды по emotion-aware
dialogue systems, распознаванию эмоций в речи, русскоязычным корпусам
эмоциональной речи, spoken dialogue architecture и промышленным
voice-pipeline решениям. Технической основой первого рабочего контура в
рамках настоящего исследования выбран стек \texttt{Silero\ VAD},
\texttt{GigaAM\ v3\ e2e\_rnnt}, \texttt{GigaAM-Emo},
\texttt{OpenAI\ GPT-5.2\ Chat} (\texttt{gpt-5.2-chat-latest}),
\texttt{Silero\ TTS} и rule-based \texttt{policy\ layer}. В качестве
основной литературной опоры для этих блоков используются обзорные и
прикладные работы {[}3; 4; 6; 7; 9{]}.

Эмпирическая проверка работы выполнена по схеме controlled
\texttt{baseline\ vs\ adaptive} сравнения. В контрольной группе
использовался неадаптивный сценарий без учета эмоционального состояния в
policy-логике, а в экспериментальной --- адаптивный сценарий, в котором
эмоция, \texttt{confidence} и стадия диалога влияли на выбор следующего
действия системы. В качестве основных метрик рассматривались доля
целевого действия, доля негативных завершений, частота эскалаций,
длительность диалога, число реплик до исхода и технические показатели
устойчивости пайплайна. Такой дизайн согласуется с рекомендациями
литературы по emotion-aware task-oriented dialogue и industrial voice
pipeline, где подчеркивается необходимость совместной оценки
конверсионных, поведенческих и технических показателей {[}1; 7; 8{]}.

Структура работы соответствует цели и задачам исследования. В первой
главе выполняется аналитический обзор современных подходов к
emotion-aware dialogue, SER, русскоязычным корпусам и архитектурам
речевых систем. Во второй главе излагается авторская методика
адаптивного автообзвона, включая архитектуру, рабочий набор
эмоциональных состояний, правила \texttt{policy\ layer}, ограничения и
требования к логированию. В третьей главе описываются прототип,
экспериментный контур, метрики и протокол A/B-проверки предложенной
методики. Заключение содержит итоговые выводы, ограничения и направления
дальнейшего развития.

\section{1 Аналитический
обзор}\label{ux430ux43dux430ux43bux438ux442ux438ux447ux435ux441ux43aux438ux439-ux43eux431ux437ux43eux440}

\subsection{1.1 Актуальность emotion-aware подходов в разговорных
системах}\label{ux430ux43aux442ux443ux430ux43bux44cux43dux43eux441ux442ux44c-emotion-aware-ux43fux43eux434ux445ux43eux434ux43eux432-ux432-ux440ux430ux437ux433ux43eux432ux43eux440ux43dux44bux445-ux441ux438ux441ux442ux435ux43cux430ux445}

Развитие больших языковых моделей и современных речевых технологий
привело к существенному росту интереса к разговорным системам, способным
не только корректно распознавать смысл пользовательской реплики, но и
учитывать эмоциональный контекст взаимодействия. Для задач исходящего
автообзвона данная тенденция имеет особое значение, поскольку
результативность диалога зависит не только от содержания предложения, но
и от того, насколько система своевременно реагирует на признаки
заинтересованности, раздражения, усталости или конфликтности со стороны
собеседника.

В современной литературе можно выделить устойчивый переход от
статических сценариев разговора к emotion-aware системам, в которых
эмоциональное состояние пользователя рассматривается как фактор выбора
следующего действия диалогового менеджера. Работы Feng и соавторов, Zhu
и соавторов, а также ряд более поздних исследований показывают, что
эмоция может быть встроена не только в генерацию ответа, но и в контур
task-oriented dialogue как параметр управления пониманием, политикой и
формированием следующей реплики. Такой вывод особенно важен для
настоящей работы, поскольку он позволяет рассматривать эмоциональную
адаптацию не как декоративное повышение эмпатичности, а как управляемое
изменение сценарной тактики {[}1; 2; 10{]}.

Вместе с тем значительная часть существующих исследований ориентирована
либо на текстовый диалог, либо на открытые conversational agents, где
прикладная цель разговора формализована слабо. Для задач исходящего
обзвона этого недостаточно. В прикладном контуре требуется не только
реагировать на эмоциональный тон речи, но и обеспечивать достижение
целевого действия, контролировать длительность диалога, ограничивать
риск конфликта и сохранять возможность объяснения принятых решений.
Следовательно, для предметной области автообзвона необходимо не просто
перенести отдельные идеи empathetic dialogue, а сформировать
воспроизводимую методику, учитывающую эмоции на уровне policy-логики и
ограничений сценария.

\subsection{1.2 Emotion-aware dialogue systems: переход от стилизации к
policy-уровню}\label{emotion-aware-dialogue-systems-ux43fux435ux440ux435ux445ux43eux434-ux43eux442-ux441ux442ux438ux43bux438ux437ux430ux446ux438ux438-ux43a-policy-ux443ux440ux43eux432ux43dux44e}

Существенная часть работ по emotion-aware dialogue показывает, что
одного лишь эмоционально окрашенного текста ответа недостаточно для
устойчивого улучшения взаимодействия. Более перспективным направлением
является включение эмоции пользователя в процедуру выбора следующего
действия системы. Именно это отличает современные policy-oriented
подходы от ранних систем, где эмоция использовалась преимущественно как
модификатор тона реплики.

Исследования в области task-oriented dialogue подтверждают, что эмоция
может рассматриваться как отдельный сигнал среды, влияющий на
планирование следующего шага. В таких системах задача состоит не просто
в том, чтобы сделать ответ «доброжелательнее», а в том, чтобы
скорректировать траекторию разговора: сократить избыточное давление,
предложить альтернативу, углубить оффер или прервать автоматическое
убеждение. Для ВКР это особенно важно, поскольку исходящий обзвон
является жестко регламентированным сценарием, где каждый шаг диалога
должен быть связан с определенной бизнес-целью {[}1; 2{]}.

Литература также показывает, что emotion-aware policy может быть
реализована разными способами: через reinforcement learning, через
управление набором допустимых действий, через planner-модули с LLM
priors либо через гибридные архитектуры, совмещающие детерминированные
правила и генеративный ответ. Для прикладного автообзвона наиболее
реалистичным выглядит именно гибридный вариант. В нем критические
решения --- например, переход к деэскалации, перевод на оператора или
возврат к безопасному нейтральному сценарию --- должны контролироваться
явными правилами, тогда как языковая модель реализует уже выбранную
тактику в естественном русском языке {[}2; 5; 10{]}.

Таким образом, аналитический обзор показывает, что современное состояние
исследований поддерживает центральную гипотезу настоящей работы: эмоция
должна влиять не только на стиль ответа, но и на политику выбора
следующего действия системы. Однако готовой воспроизводимой схемы именно
для русскоязычного исходящего автообзвона в литературе в явном виде не
представлено. Это и образует содержательную нишу для авторской методики.

\subsection{1.3 Speech-first системы и архитектуры spoken
dialogue}\label{speech-first-ux441ux438ux441ux442ux435ux43cux44b-ux438-ux430ux440ux445ux438ux442ux435ux43aux442ux443ux440ux44b-spoken-dialogue}

Для темы адаптивного автообзвона принципиально важно, что взаимодействие
в системе происходит не в текстовой, а в речевой форме. Следовательно,
эмоция должна анализироваться не после завершения диалога и не как
вспомогательная аннотация корпуса, а как оперативный сигнал в текущем
речевом контуре.

Современные исследования spoken dialogue systems подтверждают, что
эмоциональная информация может быть встроена непосредственно в
архитектуру голосового взаимодействия. В ряде работ используются speech
encoders, эмоциональные эмбеддинги, специальные механизмы выравнивания
паралингвистических признаков и каскадные либо полуинтегрированные схемы
генерации речевого ответа. В частности, работы в духе E-chat, EmoSDS и
сходных систем демонстрируют, что эмоциональный контекст способен
улучшать выбор ответа в речевом диалоге, если он учитывается до этапа
генерации, а не только после него {[}4; 5; 11{]}.

Для настоящей ВКР особенно важен вывод о практической пригодности
каскадной схемы \texttt{ASR\ +\ SER\ +\ LLM\ +\ TTS}. Хотя в литературе
обсуждаются и более сложные end-to-end speech-native системы, каскадная
архитектура сохраняет ряд преимуществ: модульность, воспроизводимость,
возможность поэтапного тестирования, раздельный анализ ошибок и
относительную инженерную доступность. В условиях магистерской работы это
делает ее более подходящей, чем полностью интегрированные speech
foundation models, которые сложнее в настройке, труднее объясняются и
хуже подходят для controlled A/B-эксперимента {[}9; 12{]}.

Наряду с преимуществами каскада литература указывает и на его
ограничения. Ключевыми проблемами являются накопление ошибок при
переходе между модулями, дополнительная задержка и частичная потеря
информации при преобразовании речи в текст. Для автообзвона особенно
критична задержка ответа, поскольку при телефонном взаимодействии
пользователь чувствителен к искусственным паузам, сбоям turn-taking и
неестественным реакциям системы. В ряде работ по full-duplex dialogue и
spoken interaction показано, что качество восприятия диалога заметно
ухудшается уже при сравнительно небольшом росте задержки. Эти
ограничения особенно подробно обсуждаются в обзоре SpeechLM и в работах
по full-duplex dialogue agents {[}9; 13; 14{]}.

Следовательно, для предметной области ВКР обоснован не любой речевой
контур, а именно такой, который сочетает модульность и контроль с
допустимым уровнем latency. Это напрямую подводит к выбору локально
исполняемых русскоязычных модулей \texttt{ASR}, \texttt{SER} и
\texttt{TTS}, а также к необходимости явного описания ограничений
задержки в последующих главах.

\subsection{1.4 Современные подходы к распознаванию эмоций в
речи}\label{ux441ux43eux432ux440ux435ux43cux435ux43dux43dux44bux435-ux43fux43eux434ux445ux43eux434ux44b-ux43a-ux440ux430ux441ux43fux43eux437ux43dux430ux432ux430ux43dux438ux44e-ux44dux43cux43eux446ux438ux439-ux432-ux440ux435ux447ux438}

Одним из центральных технических вопросов для предлагаемой методики
является выбор модуля \texttt{SER} --- распознавания эмоционального
состояния по речи. Анализ литературы показывает, что в современных
исследованиях наиболее убедительной основой для SER являются
self-supervised speech representations. Работы, посвященные
\texttt{Wav2Vec2}, \texttt{HuBERT}, \texttt{WavLM}, \texttt{emotion2vec}
и близким архитектурам, демонстрируют их устойчивое превосходство над
более ранними сверточными и традиционными handcrafted-feature подходами
на широком спектре задач распознавания эмоций {[}6; 15; 16; 17{]}.

Важно, однако, учитывать, что универсальная сила SSL-подхода еще не
означает одинаково хорошей применимости в русскоязычном телефонном
контуре. С одной стороны, исследования действительно подтверждают, что
\texttt{HuBERT}, \texttt{WavLM} и специализированные emotion-aware
pretraining схемы являются сильной основой для SER. С другой стороны,
для прикладной ВКР по русскому автообзвону более важен не выбор «самой
модной» модели, а выбор модели, адекватной языку, домену, доступным
данным и latency-ограничениям.

Отдельная линия исследований показывает, что одного акустического канала
иногда недостаточно для уверенного распознавания эмоции, особенно в
случаях неоднозначной интонации или заметных ошибок транскрибации. В
таких ситуациях комбинация акустических и текстовых признаков, а также
мультимодальные fusion подходы могут давать дополнительный прирост
качества. Для настоящей ВКР этот вывод полезен как методическое
ограничение: основной контур может строиться на акустическом
\texttt{SER}, но в дальнейшем допустимо рассматривать текстовую
корректировку в неоднозначных случаях как одно из направлений развития
{[}18; 19; 20{]}.

В целом литературный обзор подтверждает, что современный \texttt{SER}
должен строиться не вокруг абстрактного классификатора эмоций, а вокруг
надежного back-end модуля, способного возвращать не только класс эмоции,
но и меру доверия к гипотезе. Именно этот параметр в дальнейшем
позволяет переводить теоретическую идею распознавания эмоций в
practically safe policy-механику.

\subsection{1.5 Русскоязычный SER и ограничения межъязыкового
переноса}\label{ux440ux443ux441ux441ux43aux43eux44fux437ux44bux447ux43dux44bux439-ser-ux438-ux43eux433ux440ux430ux43dux438ux447ux435ux43dux438ux44f-ux43cux435ux436ux44aux44fux437ux44bux43aux43eux432ux43eux433ux43e-ux43fux435ux440ux435ux43dux43eux441ux430}

Для настоящей работы принципиальное значение имеет тот факт, что
результаты англоязычных исследований не могут автоматически переноситься
на русский телефонный диалог. Литература по русскоязычному \texttt{SER}
показывает, что выбор корпуса, языка и типа речевого материала
существенно влияет на качество моделей. Именно поэтому в главе 2 и в
экспериментальной части необходимо обосновывать не просто наличие
\texttt{SER}, а его пригодность к русскоязычному контексту.

Наиболее значимыми ориентирами здесь являются русскоязычные корпуса
эмоций в речи, в первую очередь \texttt{DUSHA} и \texttt{RESD}.
Сопоставление этих корпусов показывает, что различия в масштабе,
качестве и разнообразии речевого материала ведут к заметным различиям в
итоговой обучаемости моделей. Из этого следует, что для первого
воспроизводимого контура ВКР предпочтительно опираться на более
репрезентативный открытый русскоязычный корпус, а дополнительные
датасеты рассматривать как средства проверки переносимости {[}6; 21{]}.

Дополнительным аргументом в пользу русско-центричного выбора является
ограниченность межъязыкового переноса для эмоций. Работы по
русскоязычному SER показывают, что модель, успешно обученная на одном
языке, может заметно терять качество при переносе на другой языковой
материал. Для ВКР это означает, что нельзя строить всю аргументацию
выбора \texttt{SER} только на англоязычных бенчмарках. Необходимо явно
фиксировать, что русскоязычный телефонный диалог требует отдельной
экспериментальной и методической валидации {[}6; 21{]}.

Именно поэтому в настоящей работе более обоснованным выглядит выбор
русскоязычно-центричного контура \texttt{GigaAM-Emo} как основного
кандидата и \texttt{HuBERT\ large}, дообученного на \texttt{DUSHA}, как
воспроизводимого baseline. Такой выбор лучше соответствует и задачам
работы, и накопленной литературной базе.

\subsection{1.6 Контакт-центры, автообзвон и промышленный контур
voice-систем}\label{ux43aux43eux43dux442ux430ux43aux442-ux446ux435ux43dux442ux440ux44b-ux430ux432ux442ux43eux43eux431ux437ux432ux43eux43d-ux438-ux43fux440ux43eux43cux44bux448ux43bux435ux43dux43dux44bux439-ux43aux43eux43dux442ux443ux440-voice-ux441ux438ux441ux442ux435ux43c}

Еще один важный вывод литературы заключается в том, что для прикладных
задач контакт-центров недостаточно рассматривать отдельно только модель
SER или отдельно языковую модель. Реальные voice-системы в
индустриальном контуре должны проектироваться как сквозной pipeline,
включающий сегментацию, транскрибацию, распознавание эмоций, policy
layer, генерацию ответа, синтез речи, логирование и механизмы эскалации.

Работы по industrial SER pipeline и affective dialogue management
подчеркивают, что качество прикладной системы определяется не только
точностью отдельного классификатора, но и тем, насколько хорошо
организованы сбор данных, разметка, балансировка классов, процедуры
валидации, правила обработки неопределенности и контрольные точки
анализа ошибок. Для исходящего автообзвона это особенно важно, поскольку
в центре внимания оказывается не теоретическая эмпатичность, а
практическая результативность и безопасность разговора {[}7; 22{]}.

Отдельное место в литературе занимают работы, посвященные сервисным
ботам, пользовательскому восприятию и trustworthiness. Они показывают,
что эмоционально-чувствительные системы могут повышать субъективное
качество взаимодействия, однако это не гарантирует автоматического роста
бизнес-метрик. Следовательно, для ВКР экспериментальная проверка должна
строиться не только на показателях «приятности» общения, но и на
конверсионных, диалоговых и технических метриках {[}1; 8{]}.

Таким образом, аналитический обзор прикладных исследований приводит к
важному для настоящей работы выводу: в автообзвоне необходимо
проектировать не просто emotion classifier, а целостную методику
адаптивного голосового взаимодействия, способную быть встроенной в
промышленный процесс и проверенной на реальных показателях.

\subsection{1.7 Ограничения существующих решений и постановка задачи
настоящей
работы}\label{ux43eux433ux440ux430ux43dux438ux447ux435ux43dux438ux44f-ux441ux443ux449ux435ux441ux442ux432ux443ux44eux449ux438ux445-ux440ux435ux448ux435ux43dux438ux439-ux438-ux43fux43eux441ux442ux430ux43dux43eux432ux43aux430-ux437ux430ux434ux430ux447ux438-ux43dux430ux441ux442ux43eux44fux449ux435ux439-ux440ux430ux431ux43eux442ux44b}

Несмотря на значительный прогресс в областях empathetic dialogue, spoken
interaction и SER, текущая литература не предлагает завершенного
решения, которое можно было бы без существенных адаптаций перенести на
задачу русскоязычного исходящего автообзвона. Выявленные ограничения
можно свести к нескольким группам.

Во-первых, значительная часть работ ориентирована либо на англоязычные
данные, либо на открытые разговорные сценарии, где цели диалога заданы
слабо и не подчинены прикладной логике контакт-центра. Во-вторых, многие
emotion-aware системы описывают влияние эмоции на стиль ответа, но не
формализуют детерминированный контур выбора следующего действия.
В-третьих, в литературе редко одновременно рассматриваются все четыре
аспекта, критичные для темы ВКР: русскоязычный речевой контур,
распознавание эмоций с учетом надежности, управляемая
\texttt{policy}-адаптация и проверка эффекта в формате
\texttt{baseline\ vs\ adaptive}.

Поэтому задача настоящей работы формулируется как разработка
воспроизводимой методики адаптивного автообзвона, в которой
эмоциональное состояние собеседника и надежность его распознавания
используются как управляющие параметры выбора следующего действия
системы. Такая постановка опирается на существующую литературу, но
одновременно закрывает выявленный пробел между современными
emotion-aware исследованиями и прикладной задачей русскоязычного
автообзвона.

\subsection{1.8 Выводы по
главе}\label{ux432ux44bux432ux43eux434ux44b-ux43fux43e-ux433ux43bux430ux432ux435}

Проведенный аналитический обзор показывает, что современные исследования
подтверждают перспективность emotion-aware подходов для разговорных
систем и указывают на необходимость учета эмоций не только при генерации
ответа, но и на уровне диалоговой политики. Показано, что для речевого
взаимодействия наиболее реалистичной и воспроизводимой схемой остается
каскадный контур \texttt{ASR\ +\ SER\ +\ LLM\ +\ TTS}, несмотря на
присущие ему ограничения по задержке и накоплению ошибок.

Установлено, что для \texttt{SER} наиболее убедительной технической
основой выступают SSL-модели, однако для русскоязычного автообзвона
необходим отдельный выбор корпусов и моделей, а прямой перенос
англоязычных результатов ограничен. Показано также, что прикладная
ценность эмоциональной адаптации должна оцениваться не только по
субъективному качеству общения, но и по бизнес- и диалоговым метрикам в
контролируемом эксперименте.

Тем самым глава 1 подводит к центральной задаче работы: разработке
формальной и воспроизводимой методики адаптивного автообзвона для
русскоязычного телефонного контура. В следующей главе эта задача
решается через описание архитектуры, рабочего набора эмоциональных
состояний, правил policy layer, ограничений безопасности и механизма
интеграции языковой модели в контур принятия решения.

\section{2 Методика адаптивного
автообзвона}\label{ux43cux435ux442ux43eux434ux438ux43aux430-ux430ux434ux430ux43fux442ux438ux432ux43dux43eux433ux43e-ux430ux432ux442ux43eux43eux431ux437ux432ux43eux43dux430}

\subsection{2.1 Требования к методике адаптивного
автообзвона}\label{ux442ux440ux435ux431ux43eux432ux430ux43dux438ux44f-ux43a-ux43cux435ux442ux43eux434ux438ux43aux435-ux430ux434ux430ux43fux442ux438ux432ux43dux43eux433ux43e-ux430ux432ux442ux43eux43eux431ux437ux432ux43eux43dux430}

Центральной задачей настоящей работы является разработка не просто
прототипа голосового ассистента, а воспроизводимой методики адаптивного
автообзвона, которая позволяет управляемо изменять сценарий исходящего
телефонного диалога на основе автоматически распознаваемого
эмоционального состояния собеседника. Такая постановка принципиально
отличает предлагаемое решение от систем, где большая языковая модель
используется как слабо ограниченный генератор ответов без явной
формализации правил выбора диалоговой тактики. Такой акцент на
детерминированном уровне policy согласуется с современными emotion-aware
task-oriented подходами, где эмоция пользователя влияет на выбор
следующего действия системы, а не только на тон текста ответа {[}1;
2{]}.

Для прикладного контура исходящего обзвона методика должна удовлетворять
одновременно нескольким требованиям. Во-первых, она должна быть
воспроизводимой, то есть другой исследователь или разработчик должен
иметь возможность восстановить логику работы системы по описанию входных
данных, составу модулей, порогам уверенности и таблице управляющих
правил. Во-вторых, она должна быть объяснимой, поскольку для контактного
центра или маркетингового агентства принципиально важно понимать, по
какой причине система изменила тон реплики, сократила оффер, предложила
альтернативу либо перевела звонок на оператора. В-третьих, методика
должна быть безопасной: эмоциональная адаптация не должна превращаться в
манипулятивное давление на абонента и не должна позволять модели
свободно принимать критические решения вне регламентированных
ограничений.

С формальной точки зрения на каждом такте диалога методика получает
набор входов, включающий аудиофрагмент реплики собеседника, текстовую
гипотезу модуля автоматического распознавания речи, эмоциональную
гипотезу и оценку ее надежности, текущее состояние сценария, историю
предыдущих тактов и набор доменных ограничений. На выходе формируется не
только текст следующей реплики, но и явное policy-решение: выбранная
тактика ответа, обновленное состояние сценария, признак перехода к
fallback-режиму, а при необходимости и признак эскалации на оператора.
Следовательно, методика должна описывать как языковую генерацию, так и
предшествующую ей процедуру выбора допустимого действия.

Корректно заданной можно считать такую методику, для которой выполнены
три условия. Первое условие состоит в наличии полного покрытия
допустимых входных состояний, когда для каждого сочетания стадии
диалога, эмоционального состояния и уровня надежности распознавания
определено следующее действие системы. Второе условие состоит в
отделении детерминированного слоя принятия решения от вероятностного
слоя языковой реализации. Третье условие заключается в наличии
логирования, позволяющего постфактум восстановить, какие сигналы
поступили в систему и почему был выбран конкретный ответ.

Таким образом, в рамках данной главы методика рассматривается как
совокупность архитектурных, алгоритмических и организационных правил,
обеспечивающих управляемую адаптацию исходящего диалога при сохранении
требований к объяснимости, безопасности и последующей экспериментальной
проверке.

С учетом поставленной цели основным результатом главы должна выступать
формализованная методика, включающая состав модулей и их функции,
описание входных и выходных данных, рабочий набор эмоциональных
состояний, правила преобразования состояния собеседника в стратегию
диалога, механизм интеграции \texttt{LLM} в контур принятия решения,
fallback-логику, ограничения применимости и требования к логированию.

\subsubsection{2.1.1 Входные и выходные данные
методики}\label{ux432ux445ux43eux434ux43dux44bux435-ux438-ux432ux44bux445ux43eux434ux43dux44bux435-ux434ux430ux43dux43dux44bux435-ux43cux435ux442ux43eux434ux438ux43aux438}

На каждом шаге взаимодействия методика получает следующие входы:

\begin{enumerate}
\def\labelenumi{\arabic{enumi})}
\tightlist
\item
  аудиосигнал реплики собеседника;
\item
  результаты сегментации речевой активности;
\item
  текстовую гипотезу, полученную модулем \texttt{ASR};
\item
  оценку эмоционального состояния, полученную модулем \texttt{SER};
\item
  значение \texttt{confidence} по эмоциональной гипотезе;
\item
  текущее состояние диалогового сценария;
\item
  историю предыдущих реплик и решений \texttt{policy\ layer};
\item
  ограничения домена, включая допустимые формулировки, правила
  комплаенса и условия перевода на оператора.
\end{enumerate}

На выходе каждого такта методика должна формировать:

\begin{enumerate}
\def\labelenumi{\arabic{enumi})}
\tightlist
\item
  выбранную стратегию ответа;
\item
  обновленное состояние сценария;
\item
  инструкцию для \texttt{LLM} на генерацию следующей реплики;
\item
  признак необходимости эскалации на оператора либо безопасного
  завершения разговора;
\item
  запись в журнал событий для последующего анализа.
\end{enumerate}

\subsubsection{2.1.2 Критерии корректности
методики}\label{ux43aux440ux438ux442ux435ux440ux438ux438-ux43aux43eux440ux440ux435ux43aux442ux43dux43eux441ux442ux438-ux43cux435ux442ux43eux434ux438ux43aux438}

Методика считается корректно заданной, если для каждого допустимого
входного состояния определено следующее действие системы, решения
\texttt{policy\ layer} объяснимы через явные правила, низкая надежность
распознавания эмоций не приводит к опасным или трудно интерпретируемым
действиям, а роль \texttt{LLM} ограничена генерацией реплики внутри явно
заданных рамок при детерминированном управлении критическими переходами.

\subsubsection{2.1.3 Формальное представление
методики}\label{ux444ux43eux440ux43cux430ux43bux44cux43dux43eux435-ux43fux440ux435ux434ux441ux442ux430ux432ux43bux435ux43dux438ux435-ux43cux435ux442ux43eux434ux438ux43aux438}

Для придания методике строгого и воспроизводимого вида введем дискретное
описание одного такта диалога. Пусть в момент времени \(t\) состояние
системы задается кортежем

\begin{equation}
s_t = (g_t, e_t, c_t, h_t),
\end{equation}

где \(g_t\) обозначает текущую стадию сценария, \(e_t\) ---
нормализованное эмоциональное состояние собеседника, \(c_t\) ---
значение надежности эмоциональной гипотезы, \(h_t\) --- краткую историю
предшествующих тактов, включающую недавние policy-решения и признаки
устойчивости эмоции.

На вход такта поступают акустический фрагмент \(x_t\) и внутреннее
состояние сценария. Модули распознавания формируют текстовую гипотезу
\(y_t = \operatorname{ASR}(x_t)\) и первичную эмоциональную гипотезу
\((\hat e_t, \hat c_t) = \operatorname{SER}(x_t)\). После этого слой
нормализации строит рабочее состояние эмоции

\begin{equation}
e_t = N(\hat e_t, \hat c_t, h_t),
\end{equation}

где оператор \(N\) переводит детализированные классы SER в ограниченный
набор состояний методики и, при необходимости, заменяет их состоянием
неопределенности. В простейшем виде правило нормализации может быть
записано следующим образом:

\begin{equation}
e_t =
\begin{cases}
\mathrm{uncertain}, & \hat c_t < \theta_{\mathrm{low}},\\
\operatorname{map}(\hat e_t), & \hat c_t \ge \theta_{\mathrm{low}}.
\end{cases}
\end{equation}

При этом критические переключения тактики допускаются не для каждого
мгновенного значения \(e_t\), а только для устойчивого состояния
\(z_t = I(e_t, h_t)\), где оператор \(I\) реализует правило инерции.
Например, изменение на конфликтный режим разрешается либо при
\(\hat c_t \ge \theta_{\mathrm{high}}\), либо при подтверждении одного и
того же негативного состояния в нескольких последовательных тактах.

Выбор следующего действия задается policy-функцией

\begin{equation}
a_t = \pi(g_t, z_t, c_t, h_t),
\end{equation}

где \(a_t\) принадлежит конечному множеству допустимых действий:
следование базовому сценарию, углубление предложения, смягчение ответа,
уточняющий вопрос, предложение альтернативы, эскалация на оператора либо
корректное завершение разговора. После выбора действия обновление стадии
сценария может быть представлено отображением

\begin{equation}
g_{t+1} = T(g_t, a_t),
\end{equation}

где \(T\) --- функция перехода между стадиями диалога.

Языковая модель в предлагаемой методике не определяет действие
самостоятельно, а реализует уже выбранное решение. Поэтому генерация
ответа описывается как условная функция

\begin{equation}
r_t = G(y_t, a_t, C),
\end{equation}

где \(r_t\) --- текст следующей реплики, а \(C\) --- множество
ограничений сценария. К указанным ограничениям относятся краткость
ответа, запрет на давление на собеседника, запрет на выход за рамки
сценария и запрет на обещание условий, которые недоступны в реальном
контуре. Если \(z_t = \mathrm{anger}\) и
\(c_t \ge \theta_{\mathrm{high}}\), то множество ограничений \(C\)
дополнительно требует либо предложения перевода на оператора, либо
корректного завершения разговора.

Содержательно качество методики определяется не одной локальной
формулой, а совокупностью критериев. В экспериментальной части это
означает стремление к максимизации вероятности целевого исхода
\(P(\mathrm{target\_success})\) при одновременном ограничении доли
конфликтных завершений, необоснованных эскалаций и чрезмерной
длительности разговора. Таким образом, математическая формализация в
данной работе используется не ради абстрактного усложнения, а как
средство строгого описания состояний, переходов и допустимых действий
системы.

\subsection{2.2 Архитектура системы и состав
модулей}\label{ux430ux440ux445ux438ux442ux435ux43aux442ux443ux440ux430-ux441ux438ux441ux442ux435ux43cux44b-ux438-ux441ux43eux441ux442ux430ux432-ux43cux43eux434ux443ux43bux435ux439}

Предлагаемая методика реализуется в каскадной архитектуре, где обработка
речевого сигнала и принятие диалогового решения разделены на несколько
последовательных модулей. В состав контура входят детекция речевой
активности, параллельные модули \texttt{ASR} и \texttt{SER}, слой
нормализации, \texttt{policy\ layer}, построитель prompt-инструкции,
языковая модель, валидация ответа и синтез речи.

Выбор именно каскадной архитектуры опирается на работы по
emotion-sensitive spoken dialogue и обзоры speech-language систем, в
которых подчеркиваются преимущества модульности, управляемости и
раздельного анализа ошибок для прикладного голосового контура {[}4; 5;
9{]}. Общая архитектура методики представлена на рисунке 2.1.

\begin{figure}[H]
\centering
\VkrFigureArchitecture{\includegraphics[keepaspectratio,alt={Рисунок 2.1 --- Общая архитектура методики адаптивного автообзвона с разделением на речевые модули, normalization layer и policy layer}]{docx_pipeline/figures/png/ch2_architecture_pipeline.png}}
\caption{Рисунок 2.1 --- Общая архитектура методики адаптивного
автообзвона с разделением на речевые модули, normalization layer и
policy layer}
\end{figure}

Первым модулем выступает подсистема детекции речевой активности
\texttt{VAD}, которая определяет границы пользовательской реплики и тем
самым задает такт последующей обработки. Для текущего технологического
контура выбран \texttt{Silero\ VAD}, поскольку он хорошо подходит для
потокового речевого сценария и относительно прост в интеграции. После
выделения завершенного фрагмента аудио он передается в два канала
параллельной обработки.

Первый канал образует модуль автоматического распознавания речи
\texttt{ASR}. Его задача состоит в преобразовании звучащей реплики в
текстовую форму, пригодную для анализа содержания и последующей передачи
в языковую модель. В текущем цикле работы основной рабочей гипотезой
выбран \texttt{GigaAM\ v3\ e2e\_rnnt} как русскоязычно-ориентированный
вариант транскрибации, а \texttt{faster-whisper} оставлен не как
основной резерв, а как внешний baseline-компаратор. Такое решение
позволяет одновременно учитывать прикладную потребность в лучшей работе
с русской речью и сохранять воспроизводимую внешнюю точку сравнения.

Второй канал составляет модуль распознавания эмоционального состояния
\texttt{SER}. Он получает тот же аудиофрагмент и возвращает первичную
гипотезу эмоции и числовую оценку уверенности \texttt{confidence}.
Принципиально важно подчеркнуть, что SER в данной работе не используется
как самостоятельный классификатор ради отчетности. Его выход служит
управляющим сигналом для последующего выбора диалоговой тактики. Поэтому
в архитектуре после \texttt{ASR} и \texttt{SER} вводится специальный
слой нормализации, приводящий разнородные данные к единому формату
принятия решения.

\texttt{Normalization\ layer} выполняет несколько функций. Во-первых, он
агрегирует первичные классы эмоций в ограниченный рабочий набор
состояний, пригодный для прикладного сценария автообзвона. Во-вторых, он
учитывает значение \texttt{confidence}, переводя слабые либо
противоречивые гипотезы в безопасное состояние неопределенности.
В-третьих, он может учитывать историю ближайших тактов, чтобы устранять
кратковременные шумовые всплески и не допускать избыточно частого
переключения тактики.

Ключевым элементом архитектуры является \texttt{policy\ layer}. В
отличие от свободной генерации текста, этот слой реализует
детерминированное отображение emotion + confidence + dialogue stage →
action. Иначе говоря, именно он отвечает за выбор того, что система
должна сделать на следующем шаге: продолжить базовый сценарий, углубить
предложение, смягчить коммуникацию, задать уточняющий вопрос, выполнить
fallback либо инициировать эскалацию. Тем самым методика сохраняет
управляемость и пригодность к аудиту. Такой способ преобразования
эмоционального сигнала в управляемое действие концептуально близок к
emotion-sensitive dialogue policy и affective dialogue manager {[}2;
22{]}.

После выбора действия в работу вступает \texttt{prompt\ builder},
преобразующий policy-решение в структурированную инструкцию для языковой
модели. Далее \texttt{LLM} генерирует одну следующую реплику в
допустимых границах, а модуль \texttt{response\ validation} проверяет ее
на соответствие сценарию, ограничениям комплаенса и правилам безопасного
поведения. Только после этого итоговый ответ передается в модуль синтеза
речи \texttt{TTS}, для которого в текущем стеке основным локальным
кандидатом рассматривается \texttt{Silero\ TTS}, а облачным
сравнительным вариантом может выступать \texttt{SaluteSpeech}.

Преимуществом каскадной архитектуры является то, что она проще
объясняется, масштабируется и экспериментально проверяется в рамках
магистерской ВКР, чем полностью end-to-end голосовая система. В
предлагаемом контуре каждый модуль имеет четко определенную функцию, а
значит, ошибки распознавания, ошибки policy-логики и ошибки генерации
могут анализироваться раздельно.

\subsection{2.3 Формализация эмоционального состояния и выбор модуля
SER}\label{ux444ux43eux440ux43cux430ux43bux438ux437ux430ux446ux438ux44f-ux44dux43cux43eux446ux438ux43eux43dux430ux43bux44cux43dux43eux433ux43e-ux441ux43eux441ux442ux43eux44fux43dux438ux44f-ux438-ux432ux44bux431ux43eux440-ux43cux43eux434ux443ux43bux44f-ser}

Для задачи исходящего телефонного обзвона нецелесообразно использовать
слишком детализированный психологический словарь эмоций. Цель системы
состоит не в тонкой психометрической диагностике, а в выборе
практической стратегии дальнейшего взаимодействия. По этой причине в
работе вводится ограниченный операциональный набор состояний
собеседника: нейтральность, заинтересованность, раздражение, злость и
неопределенность.

Состояние нейтральности соответствует случаям, когда в речи абонента не
наблюдается отчетливых сигналов сопротивления либо повышенной
вовлеченности. Для такого состояния система продолжает базовую
траекторию сценария. Состояние заинтересованности фиксирует признаки
готовности слушать предложение, задавать уточняющие вопросы либо
рассматривать следующий шаг. Раздражение отражает умеренно негативную
реакцию, связанную с неудобством, усталостью, сопротивлением формату
звонка или нежеланием тратить время. Злость соответствует выраженной
конфликтности, при которой продолжение автоматического убеждения
становится нежелательным. Наконец, состояние неопределенности
используется тогда, когда модель не дает надежной эмоциональной гипотезы
и система обязана перейти к безопасному варианту поведения.

Такой набор состояний выбран по двум причинам. Во-первых, он достаточно
мал, чтобы обеспечить стабильное отображение эмоций в конкретные
тактические действия. Во-вторых, он покрывает основные коммуникативные
режимы, критичные для сценариев исходящего обзвона: нейтральное
продолжение разговора, развитие потенциальной конверсии, деэскалация
умеренного негатива и прекращение автоматического убеждения при
конфликте.

Для выбранного основного кандидата \texttt{GigaAM-Emo} первичный выход
модели рассматривается как распределение вероятностей по четырем
классам: \texttt{angry}, \texttt{sad}, \texttt{neutral} и
\texttt{positive}. Эти классы не используются в методике напрямую, а
переводятся в рабочий набор состояний автообзвона. Класс
\texttt{neutral} отображается в нейтральность, \texttt{positive} --- в
заинтересованность, \texttt{sad} --- в умеренный негативный режим, для
которого применяется тактика раздражения и снижения давления, а
\texttt{angry} --- в состояние злости. При отсутствии устойчивой
гипотезы либо при низкой уверенности используется состояние
неопределенности.

Особое значение в методике имеет параметр \texttt{confidence}. В ряде
прикладных работ уверенность классификатора рассматривается как
вспомогательный показатель, однако для адаптивного автообзвона такой
подход недостаточен. Если система будет одинаково реагировать на слабую
и на надежную эмоциональную гипотезу, она станет либо слишком
чувствительной к шуму, либо слишком рискованной с точки зрения
конфликтных ситуаций. Поэтому \texttt{confidence} включается в
формальное описание методики как обязательный управляющий параметр.
Необходимость такого решения дополнительно подтверждается
исследованиями, где качество SER заметно зависит от неоднозначности
сигнала, ошибок ASR и мультимодального контекста {[}15; 18{]}.

В первом контуре методики вводятся два порога надежности эмоциональной
гипотезы: \(\theta_{\mathrm{low}}\) и \(\theta_{\mathrm{high}}\). В
рамках пилотного эксперимента они используются как настраиваемые
параметры policy layer, а не как универсальные константы, переносимые
между всеми кампаниями без калибровки. Если
\(c_t < \theta_{\mathrm{low}}\), то первичный эмоциональный вывод
считается ненадежным, а система переводит такт в состояние
неопределенности. Если значение находится в диапазоне от
\(\theta_{\mathrm{low}}\) до \(\theta_{\mathrm{high}}\), допускается
только мягкая адаптация без критических переходов. И только при
\(c_t \ge \theta_{\mathrm{high}}\) разрешается применение полного набора
policy-действий. Отдельное правило задается для состояния злости: при
достаточно высокой уверенности, например выше 0,9, система должна
отказаться от автоматического убеждения и перейти к эскалации либо к
корректному завершению разговора. В дальнейшем точные значения порогов
должны уточняться по данным конкретной кампании, поскольку качество
эмоциональной классификации зависит от оффера, состава клиентской базы и
характеристик телефонного канала.

С точки зрения выбора технической основы для \texttt{SER} в
рассматриваемой реализации целесообразно ориентироваться на
русскоязычно-центричный контур, в котором основным кандидатом выступает
\texttt{GigaAM-Emo}, а воспроизводимый baseline может быть построен на
\texttt{HuBERT\ large}, дообученном на датасете \texttt{DUSHA}. Такое
решение лучше согласуется с прикладной задачей русского телефонного
диалога, чем абстрактный выбор универсальной SSL-модели без привязки к
русскому корпусу эмоциональной речи. При этом в тексте ВКР важно
подчеркнуть, что конкретная модель \texttt{SER} является реализационным
уровнем методики, тогда как центральным авторским результатом остается
именно способ использования эмоциональной гипотезы и ее надежности в
диалоговом policy-контуре. Такая аргументация опирается на русскоязычные
обзоры и экспериментальные работы по корпусам \texttt{DUSHA} и
\texttt{RESD}, а также на общую линию исследований по SSL-архитектурам
для SER {[}6; 16; 21{]}.

\subsection{2.4 Алгоритм адаптации сценария и policy
layer}\label{ux430ux43bux433ux43eux440ux438ux442ux43c-ux430ux434ux430ux43fux442ux430ux446ux438ux438-ux441ux446ux435ux43dux430ux440ux438ux44f-ux438-policy-layer}

Основу предлагаемой методики составляет правило преобразования тройки
\texttt{эмоциональное\ состояние,\ confidence,\ стадия\ сценария} в
следующее действие системы. Такое действие включает не только
тематическое содержание ответа, но и его тактические параметры: степень
эмпатии, краткость, глубину предложения, допустимость уточняющих
вопросов, необходимость альтернативного шага и условия остановки
автоматического взаимодействия.

Для прикладного сценария исходящего обзвона целесообразно выделить шесть
стадий диалога: установление контакта, краткое представление цели
звонка, выявление интереса и первичную квалификацию, краткую
конкретизацию ценности, обработку возражения или сомнения, а также
фиксацию следующего шага либо завершение разговора. Наличие стадии
обязательно, поскольку одна и та же эмоция должна по-разному
интерпретироваться в зависимости от того, на каком этапе находится
взаимодействие. Например, заинтересованность на стадии первичной
квалификации должна вести к углублению оффера, тогда как
заинтересованность на завершающей стадии должна переводиться в краткое
закрепление следующего действия, а не в повторное развертывание
предложения. Такой stage-aware принцип соответствует современным
emotion-aware моделям task-oriented и spoken dialogue, где политика
определяется не только локальной эмоцией, но и местом реплики в
структуре сценария {[}1; 2; 5{]}.

Базовая логика policy layer может быть описана следующими правилами. При
нейтральном состоянии система следует базовому сценарию, сохраняет
краткость и избегает избыточной детализации. При заинтересованности она
усиливает информативность реплики, конкретизирует ценность предложения и
предлагает следующий шаг конверсии. При раздражении система снижает
темп, признает неудобство разговора и предлагает либо короткий формат
продолжения, либо альтернативу, например перенос контакта. При злости
она отказывается от persuasion-скрипта и инициирует деэскалацию, перевод
на оператора или корректное завершение. При неопределенности методика
предписывает безопасный нейтральный режим с мягким уточняющим вопросом
либо безусловным возвратом к базовому сценарию.

Существенным элементом алгоритма является правило инерции. Без него
система может начать резко менять тактику из-за кратковременного
эмоционального шума, ошибки классификатора или нестабильности канала
связи. Поэтому в методике предлагается допускать критические
переключения только в двух случаях: либо состояние подтверждается в
нескольких последовательных тактах, либо его надежность превышает
высокий порог. Это позволяет одновременно сохранить чувствительность к
действительно важным изменениям и снизить вероятность ложной адаптации.
Подобная осторожная логика особенно важна для прикладных voice систем,
где шум канала и ошибки промежуточных модулей могут приводить к ложным
эмоциональным срабатываниям {[}7; 18{]}.

Логика одного такта обработки может быть представлена в следующей
последовательности. После завершения реплики абонента система выделяет
аудиофрагмент, одновременно подает его в \texttt{ASR} и \texttt{SER},
нормализует эмоциональное состояние, соотносит его с \texttt{confidence}
и текущей стадией сценария, выбирает policy-действие, строит
структурированный prompt для языковой модели, генерирует реплику,
проверяет ее на допустимость, озвучивает через \texttt{TTS} и записывает
все промежуточные результаты в журнал. Тем самым каждый шаг адаптации
оказывается наблюдаемым и пригодным для анализа.

Важно подчеркнуть, что воспроизводимость методики обеспечивается не
фактом использования одной и той же LLM, а наличием явных точек
детерминации. К ним относятся фиксированный набор состояний, пороги
уверенности, таблица правил, условия fallback-перехода и логирование
итогового policy-действия. Благодаря этому экспериментальная часть
работы может проверять именно эффект методики, а не случайную
вариативность генерации текста.

Для исключения содержательной неэквивалентности между группами
целесообразно явно зафиксировать различие между baseline- и
adaptive-сценарием именно на уровне тактики, а не на уровне воронки
звонка. В baseline-варианте система проходит те же стадии разговора,
однако действует по фиксированной последовательности реплик и не
использует выход \texttt{SER} в \texttt{policy\ layer}. В
adaptive-варианте цель звонка, допустимые исходы и стадийная структура
сохраняются, но конкретный следующий шаг меняется в зависимости от
эмоции и значения \texttt{confidence}. Иначе говоря, adaptive-контур не
подменяет сценарий целиком, а модифицирует его внутри заранее заданной
рамки.

С практической точки зрения стадийная логика двух режимов может быть
описана следующим образом. На стадии установления контакта
baseline-подход проверяет возможность разговора по стандартному шаблону,
тогда как adaptive-подход при раздражении сокращает вступление, а при
злости высокой уверенности быстро переходит к корректному завершению
либо к переводу на оператора. На стадии краткого представления цели
звонка baseline ограничивается одной и той же формулой оффера, а
adaptive допускает более предметную и короткую подачу при
заинтересованности либо смягченную формулировку с акцентом на
добровольность продолжения при негативной реакции.

На стадии выявления интереса и первичной квалификации baseline
использует фиксированный вопрос о релевантности предложения, тогда как
adaptive при заинтересованности углубляет квалификацию, при
неопределенности задает нейтральный уточняющий вопрос, а при раздражении
быстрее переводит разговор к короткой развилке. На стадии краткой
конкретизации ценности baseline сообщает ограниченный набор аргументов в
неизменном виде, в то время как adaptive при заинтересованности
добавляет релевантную конкретику и предлагает следующий шаг, а при
раздражении, напротив, не расширяет оффер и быстрее переходит к
альтернативе. Наконец, на стадиях обработки возражения и фиксации
следующего шага baseline опирается на заранее подготовленные шаблоны без
учета эмоции, тогда как adaptive либо снижает давление и предлагает
перенос, либо закрепляет согласованный следующий шаг, либо прекращает
persuasion-скрипт при признаках конфликта. Стадийная логика baseline- и
adaptive-сценариев показана на рисунке 2.2.

\begin{figure}[H]
\centering
\VkrFigureStagePolicy{\includegraphics[keepaspectratio,alt={Рисунок 2.2 --- Стадийная логика baseline- и adaptive-сценариев при сохранении общей воронки звонка}]{docx_pipeline/figures/png/ch2_stage_policy_map.png}}
\caption{Рисунок 2.2 --- Стадийная логика baseline- и adaptive-сценариев
при сохранении общей воронки звонка}
\end{figure}

\subsection{2.5 Интеграция LLM в контур принятия
решения}\label{ux438ux43dux442ux435ux433ux440ux430ux446ux438ux44f-llm-ux432-ux43aux43eux43dux442ux443ux440-ux43fux440ux438ux43dux44fux442ux438ux44f-ux440ux435ux448ux435ux43dux438ux44f}

Использование большой языковой модели в рамках предлагаемой методики
имеет строго ограниченную функцию. Для текущего рабочего контура в
качестве основной \texttt{LLM} фиксируется
\texttt{OpenAI\ GPT-5.2\ Chat} (\texttt{gpt-5.2-chat-latest}).
\texttt{LLM} не должна выступать автономным центром управления диалогом,
самостоятельно определяющим, как реагировать на эмоцию собеседника и
когда завершать сценарий. Ее роль состоит в языковой реализации уже
выбранной тактики. Иначе говоря, policy layer отвечает на вопрос о том,
какое действие следует выполнить, а языковая модель отвечает на вопрос о
том, как это действие сформулировать в естественном русском языке.

Такое разграничение необходимо по нескольким причинам. Во-первых, оно
повышает объяснимость системы: критические переходы можно объяснить не
внутренним состоянием нейросети, а явным правилом. Во-вторых, оно
уменьшает риск непредсказуемого поведения в чувствительных ситуациях,
например при конфликтной реакции абонента. В-третьих, оно делает
возможным воспроизводимый эксперимент, в котором исследуется влияние
policy layer, а не произвольных вариаций свободного диалогового
поведения модели. Такое разделение между policy и generation согласуется
с подходами, где эмоциональный контекст участвует в управлении диалогом,
но не подменяет его свободной генерацией {[}1; 2; 5{]}.

Для передачи решения policy layer языковой модели используется
структурированный prompt. В его состав должны входить следующие блоки:
роль системы, стадия сценария, цель текущего шага, краткая история
диалога, последняя реплика пользователя в текстовой форме,
нормализованное эмоциональное состояние, значение \texttt{confidence},
выбранная тактика и набор ограничений. К ограничениям относятся запрет
на выход за рамки сценария, запрет на обещание недоступных условий,
ограничение на длину ответа и специальное правило деэскалации при
высокой вероятности злости.

В зависимости от нормализованного состояния структура prompt остается
общей, но меняются локальные инструкции к ответу. При нейтральности
формируется спокойная краткая реплика без эмоционального усиления. При
заинтересованности уточняется ценность предложения и следующий шаг. При
раздражении ответ признает неудобство и предлагает короткий или
альтернативный вариант. При злости система прекращает убеждение и
переходит к деэскалации. При неопределенном состоянии задается мягкий
уточняющий вопрос либо используется базовый сценарий.

Таким образом, \texttt{LLM} в предлагаемой архитектуре выступает не
источником правил, а языковым исполнителем policy-решения. Такое
понимание роли языковой модели согласуется с целью работы: построить
методику, а не заменить логику диалогового управления одной генеративной
подсистемой.

С практической точки зрения это означает, что в prompt должны
передаваться не произвольные пожелания к генерации, а формально заданные
параметры текущего такта: стадия сценария, цель текущего шага, краткая
история диалога, текст последней реплики абонента, нормализованное
эмоциональное состояние, значение \texttt{confidence}, выбранное
policy-действие и набор ограничений. Последний набор должен включать
запрет на давление на собеседника, запрет на выход за рамки сценария,
запрет на обещание недоступных условий, обязательную деэскалацию при
высокой уверенности в состоянии злости и требование краткости ответа,
достаточной для естественного озвучивания. Тем самым в текст главы
включается не техническая памятка, а содержательное описание
ограниченного класса допустимых инструкций для \texttt{LLM}.

\subsection{2.6 Fallback-логика, требования к логированию и ограничения
методики}\label{fallback-ux43bux43eux433ux438ux43aux430-ux442ux440ux435ux431ux43eux432ux430ux43dux438ux44f-ux43a-ux43bux43eux433ux438ux440ux43eux432ux430ux43dux438ux44e-ux438-ux43eux433ux440ux430ux43dux438ux447ux435ux43dux438ux44f-ux43cux435ux442ux43eux434ux438ux43aux438}

Для практического применения в контуре исходящего обзвона методика
должна содержать не только основной сценарий адаптации, но и набор
fallback-механизмов, которые предотвращают опасные либо плохо
интерпретируемые действия системы. К основаниям перехода в fallback
относятся низкая уверенность \texttt{SER}, существенные ошибки
\texttt{ASR}, невозможность надежно интерпретировать последнюю реплику
собеседника, непрохождение ответа языковой модели через модуль
валидации, повторяющийся конфликтный паттерн и чрезмерная техническая
задержка одного такта диалога.

Fallback-реакции должны быть многоуровневыми. Наиболее мягкий вариант
состоит в возврате к нейтральному сценарию без эмоционально окрашенной
адаптации. Следующим вариантом является уточняющий вопрос, позволяющий
безопасно получить более ясный сигнал о намерении собеседника. Далее
могут использоваться предложение продолжить позже, перевод на оператора
либо корректное завершение разговора. Такой набор действий важен не
только с инженерной точки зрения, но и как обязательный элемент
исследовательской методики, поскольку именно fallback-логика определяет
поведение системы в неопределенных и предельных случаях.

Не менее важным требованием является журналирование всех значимых
событий диалога. Для каждого такта целесообразно сохранять идентификатор
звонка, временную метку, длительность пользовательской реплики, текст
\texttt{ASR}, первичный выход \texttt{SER}, значение
\texttt{confidence}, нормализованное состояние, текущую стадию сценария,
выбранное policy-действие, итоговую реплику системы, факт fallback либо
эскалации и финальный исход разговора. Такое логирование необходимо как
для инженерного сопровождения прототипа, так и для последующей научной
оценки предлагаемой методики в A/B-эксперименте.

Вместе с тем для главы 3 важно заранее разделить журналирование на
уровень всего звонка и уровень отдельного такта. На уровне звонка должны
храниться идентификатор контакта, экспериментальная группа, временное
окно, итоговый исход разговора, признак достижения целевого действия,
суммарное число тактов, общая длительность, факт перевода на оператора и
технический статус завершения. На уровне такта должны фиксироваться
\texttt{turn\_id}, временные метки, стадия сценария до и после ответа,
длительность пользовательской реплики, текст \texttt{ASR}, первичный
класс \texttt{SER}, значение \texttt{confidence}, нормализованное
состояние, выбранное \texttt{policy}-действие, признак
\texttt{fallback}, текст итоговой системной реплики и длительность
генерации и синтеза. Такое двухуровневое разделение позволяет отдельно
анализировать итог звонка и внутреннюю динамику принятия решений.

Наличие детализированного журнала позволяет решать несколько задач.
Во-первых, оно дает возможность анализировать ложные адаптации и ошибки
нормализации эмоций. Во-вторых, оно обеспечивает сравнимость групп
baseline и adaptive в экспериментальной главе. В-третьих, оно открывает
возможность для абляционного анализа порогов \texttt{confidence},
отдельных policy-правил и условий эскалации. Следовательно, логирование
должно рассматриваться не как вспомогательная инженерная опция, а как
обязательное условие воспроизводимости и научной валидности результатов.
Такой подход соответствует работам по industrial pipeline для call
center, где логирование, разметка и контроль ошибок являются неотделимой
частью экспериментального контура {[}7; 9{]}. Логика одного такта
диалога с веткой fallback показана на рисунке 2.3.

\begin{figure}[H]
\centering
\VkrFigureTurnCycle{\includegraphics[keepaspectratio,alt={Рисунок 2.3 --- Логика одного такта диалога с веткой fallback и обязательным журналированием событий}]{docx_pipeline/figures/png/ch2_turn_cycle.png}}
\caption{Рисунок 2.3 --- Логика одного такта диалога с веткой fallback и
обязательным журналированием событий}
\end{figure}

Минимальная структура журнала одного такта включает следующие поля:

\begin{enumerate}
\def\labelenumi{\arabic{enumi})}
\tightlist
\item
  идентификатор звонка;
\item
  временную метку;
\item
  длительность пользовательской реплики;
\item
  текст \texttt{ASR};
\item
  первичный выход \texttt{SER};
\item
  значение \texttt{confidence};
\item
  нормализованное состояние;
\item
  стадию сценария;
\item
  выбранное \texttt{policy}-действие;
\item
  итоговую реплику системы;
\item
  факт \texttt{fallback} либо эскалации;
\item
  финальный исход разговора.
\end{enumerate}

Итог разговора при этом должен разметаться по единой таксономии исходов:
\texttt{target\_\allowbreak{}success}, \texttt{qualified\_\allowbreak{}no\_\allowbreak{}commitment},
\texttt{soft\_\allowbreak{}refusal}, \texttt{hard\_\allowbreak{}refusal},
\texttt{human\_\allowbreak{}escalation}, \texttt{technical\_\allowbreak{}failure} и
\texttt{no\_\allowbreak{}contact}. Под \texttt{target\_\allowbreak{}success} в текущем пилотном
сценарии следует понимать явно подтвержденное согласие абонента на
следующий содержательный шаг, например на перевод на оператора,
консультацию, демонстрацию либо повторный звонок в конкретный
согласованный интервал. Такая трактовка необходима, чтобы конверсионная
метрика в главе 3 считалась по одному и тому же правилу в обеих
экспериментальных группах.

При этом в главе необходимо явно зафиксировать и ограничения
предлагаемой методики. Прежде всего, точность \texttt{SER} может
существенно снижаться на шумных каналах связи и вне домена обучающих
данных, что особенно важно для русскоязычного обзвона с неоднородным
качеством телефонии. Кроме того, каскад
\texttt{ASR\ +\ SER\ +\ LLM\ +\ TTS} подвержен накоплению ошибок:
неточная транскрибация может исказить контекст реплики, а ошибка
эмоциональной классификации может привести к выбору неоптимальной
тактики. Следует учитывать и то, что эмоциональное состояние человека не
исчерпывается дискретными классами, а потому предлагаемая нормализация
сознательно упрощает реальную картину ради управляемости прикладного
сценария.

Наконец, методика применима не для любых разговорных систем, а прежде
всего для регламентированных сценариев исходящего обзвона, где можно
выделить стадии диалога, допустимые действия, ограничения комплаенса и
условия эскалации. Именно в таких условиях эмоциональная адаптация
способна стать контролируемым и экспериментально проверяемым механизмом
повышения качества взаимодействия, а не неконтролируемой вариацией
свободного генеративного диалога.

\subsubsection{2.6.1 Базовая таблица правил
адаптации}\label{ux431ux430ux437ux43eux432ux430ux44f-ux442ux430ux431ux43bux438ux446ux430-ux43fux440ux430ux432ux438ux43b-ux430ux434ux430ux43fux442ux430ux446ux438ux438}

В целях однозначной репликации методики базовые правила адаптации удобно
зафиксировать в табличной форме; они приведены в таблице 3.

\VkrTableCaption{Таблица 3 --- Базовые правила адаптации сценария по распознанному эмоциональному состоянию собеседника}

{\def\LTcaptype{none} % do not increment counter
\begin{longtable}[]{|%
>{\RaggedRight\arraybackslash}p{(\linewidth - 15\tabcolsep) * \real{0.2500}}|
  >{\RaggedRight\arraybackslash}p{(\linewidth - 15\tabcolsep) * \real{0.1875}}|
  >{\RaggedRight\arraybackslash}p{(\linewidth - 15\tabcolsep) * \real{0.1875}}|
  >{\RaggedRight\arraybackslash}p{(\linewidth - 15\tabcolsep) * \real{0.1875}}|
  >{\RaggedRight\arraybackslash}p{(\linewidth - 15\tabcolsep) * \real{0.1875}}|}
\hline
\begin{minipage}[b]{\linewidth}\RaggedRight
Состояние
\end{minipage} & \begin{minipage}[b]{\linewidth}\RaggedRight
Условия
\end{minipage} & \begin{minipage}[b]{\linewidth}\RaggedRight
Тактика системы
\end{minipage} & \begin{minipage}[b]{\linewidth}\RaggedRight
Ограничения
\end{minipage} & \begin{minipage}[b]{\linewidth}\RaggedRight
Ожидаемый эффект
\end{minipage} \\
\hline
\endfirsthead
\hline\multicolumn{5}{|l|}{\small Продолжение таблицы 3}\\
\hline
\begin{minipage}[b]{\linewidth}\RaggedRight
Состояние
\end{minipage} & \begin{minipage}[b]{\linewidth}\RaggedRight
Условия
\end{minipage} & \begin{minipage}[b]{\linewidth}\RaggedRight
Тактика системы
\end{minipage} & \begin{minipage}[b]{\linewidth}\RaggedRight
Ограничения
\end{minipage} & \begin{minipage}[b]{\linewidth}\RaggedRight
Ожидаемый эффект
\end{minipage} \\
\hline
\endhead
\hline
\endfoot
\hline
\endlastfoot
Нейтральность & \texttt{confidence} достаточен, негативных сигналов нет.
& Следовать базовому сценарию, уточнять потребности, сохранять
краткость. & Не перегружать деталями. & Сохранение управляемого темпа
диалога. \\
Заинтересованность & Есть признаки вовлеченности и готовности слушать. &
Углубить предложение, добавить конкретику, предложить следующий шаг. &
Не переходить к давлению и манипуляции. & Рост вероятности целевого
действия. \\
Раздражение & Негативные сигналы умеренные, контакт еще сохраняется. &
Снизить темп, подтвердить неудобство, предложить короткий формат или
альтернативу. & Не спорить, не повторять оффер в прежней форме. &
Снижение вероятности досрочного завершения. \\
Злость & Высокий негатив, резкая интонация, hostile cues. & Прекратить
убеждение, предложить перевод на оператора или корректное завершение. &
Запрещено продолжать убеждающий скрипт. & Снижение риска эскалации
конфликта. \\
Неопределенность & Низкая уверенность \texttt{SER} или конфликт
сигналов. & Использовать базовый сценарий и короткий уточняющий вопрос.
& Не выполнять критические переходы. & Сохранение устойчивости и
объяснимости. \\
\end{longtable}
}

\subsubsection{2.6.2 Псевдокод одного
такта}\label{ux43fux441ux435ux432ux434ux43eux43aux43eux434-ux43eux434ux43dux43eux433ux43e-ux442ux430ux43aux442ux430}

Для обеспечения воспроизводимости логика одного такта записана в виде
псевдокода, приведенного в листинге 1.

\Needspace{0.68\textheight}
\begin{lstlisting}[caption={Псевдокод одного такта адаптивного автообзвона},label={lst:adaptive-turn}]
asr_text <- ASR(audio_chunk)
ser_emotion, ser_confidence <- SER(audio_chunk)
emotion <- normalize_emotion(ser_emotion, ser_confidence)

if ser_confidence < theta_low:
    emotion <- UNCERTAIN

action <- select_policy(
    emotion, ser_confidence, scenario_stage, history)

if emotion == ANGER and ser_confidence > 0.9:
    action <- ESCALATE_TO_HUMAN

prompt <- build_prompt(scenario_stage, asr_text, emotion, action)
response <- validate_response(LLM(prompt), compliance_rules, action)

if response is invalid:
    response <- fallback_template(action)

tts_audio <- TTS(response)

log_event(
    asr_text, ser_emotion, ser_confidence,
    emotion, action, response
)
return tts_audio, update_state(dialogue_state, action)
\end{lstlisting}

\subsection{2.7 Выводы по
главе}\label{ux432ux44bux432ux43eux434ux44b-ux43fux43e-ux433ux43bux430ux432ux435-1}

В данной главе сформулирована методика адаптивного автообзвона,
основанная на разделении функций между речевыми модулями обработки
сигнала, детерминированным policy-слоем и языковой моделью как средством
генерации реплики. Показано, что центральным элементом методики является
использование не только класса эмоции, но и надежности его
распознавания, а также стадии сценария как управляющих параметров выбора
следующего действия системы.

Предложенная методика задает ограниченный рабочий набор эмоциональных
состояний, правила их нормализации, пороговую интерпретацию
\texttt{confidence}, базовую таблицу адаптивных тактик, условия
деэскалации и требования к fallback-логике. За счет этого обеспечиваются
объяснимость и воспроизводимость поведения системы, что принципиально
важно для последующей экспериментальной проверки.

Следующим шагом исследования выступает описание прототипа и протокола
экспериментальной верификации предложенной методики на основе
A/B-сравнения неадаптивного и адаптивного сценариев.

\section{3 Экспериментальная
часть}\label{ux44dux43aux441ux43fux435ux440ux438ux43cux435ux43dux442ux430ux43bux44cux43dux430ux44f-ux447ux430ux441ux442ux44c}

\subsection{3.1 Цель экспериментальной проверки и логика
главы}\label{ux446ux435ux43bux44c-ux44dux43aux441ux43fux435ux440ux438ux43cux435ux43dux442ux430ux43bux44cux43dux43eux439-ux43fux440ux43eux432ux435ux440ux43aux438-ux438-ux43bux43eux433ux438ux43aux430-ux433ux43bux430ux432ux44b}

Экспериментальная часть настоящей работы направлена на проверку того,
что предложенная методика адаптивного автообзвона обладает не только
теоретической обоснованностью, но и прикладной эффективностью по
сравнению с базовым неадаптивным сценарием. Если во второй главе
методика была формализована как система правил, модулей и ограничений,
то в данной главе она рассматривается с точки зрения экспериментальной
верификации в контролируемом прикладном контуре.

Цель эксперимента состоит в количественной оценке влияния эмоциональной
адаптации сценария на результативность исходящего телефонного диалога.
Соответственно, объектом экспериментальной проверки выступают
автоматизированные исходящие звонки в регламентированном сценарии,
реализованные на общем технологическом стеке
\texttt{VAD\ +\ ASR\ +\ SER\ +\ LLM\ +\ TTS}, а предметом проверки
является вклад \texttt{policy\ layer}, использующего распознанное
эмоциональное состояние и значение \texttt{confidence} для изменения
сценарной тактики ответа.

Логика главы строится таким образом, чтобы сначала зафиксировать
воспроизводимый протокол сравнения, а затем показать фактические
результаты пилотного A/B-эксперимента. В ней описываются гипотезы,
дизайн сравнения, система метрик, процедура экспертной оценки
транскриптов, статистические расчеты, интерпретация эффекта и угрозы
валидности. Такой подход соответствует прикладным работам по
emotion-aware dialogue и industrial voice pipeline, где протокол и
критерии оценки должны быть отделены от последующей интерпретации
полученных результатов {[}1; 7{]}.

\subsection{3.2 Проверяемые
гипотезы}\label{ux43fux440ux43eux432ux435ux440ux44fux435ux43cux44bux435-ux433ux438ux43fux43eux442ux435ux437ux44b}

Основная гипотеза исследования состоит в том, что адаптивный автообзвон,
использующий информацию об эмоциональном состоянии собеседника и
надежности его распознавания, позволяет повысить долю целевых исходов
разговора по сравнению с неадаптивным сценарием. Данная гипотеза
непосредственно связана с центральным тезисом работы о том, что эмоция
должна влиять не только на стиль генерации реплики, но и на
policy-решение следующего такта.

В качестве частных исследовательских гипотез в работе фиксируются
следующие положения. Во-первых, эмоциональная адаптация должна снижать
долю негативных реакций и досрочных завершений разговора. Во-вторых,
правило ранней эскалации при состоянии злости и высоком значении
\texttt{confidence} должно уменьшать риск развития конфликтного
взаимодействия. В-третьих, использование \texttt{confidence} как
управляющего параметра должно обеспечивать более устойчивое поведение
системы по сравнению с адаптацией только по классу эмоции. В-четвертых,
ожидается, что наибольший эффект адаптации проявится на стадиях
обработки возражений и предложения следующего шага, тогда как на этапе
установления контакта различия между сценариями могут быть менее
выраженными. Такая логика гипотез опирается на работы, где emotion-aware
политика оценивается не только по общему качеству ответа, но и по
структуре диалога, конфликтности и устойчивости поведения системы {[}1;
2; 8{]}.

Тем самым совокупность гипотез охватывает не только итоговую конверсию,
но и структуру диалога, качество взаимодействия и устойчивость
предложенной методики.

\subsection{3.3 Экспериментальный дизайн baseline vs
adaptive}\label{ux44dux43aux441ux43fux435ux440ux438ux43cux435ux43dux442ux430ux43bux44cux43dux44bux439-ux434ux438ux437ux430ux439ux43d-baseline-vs-adaptive}

Базовой схемой экспериментальной проверки стало A/B-сравнение двух
режимов работы системы. Состав контрольной и экспериментальной групп
приведен в таблице 4.

\VkrTableCaption{Таблица 4 --- Состав контрольной и экспериментальной групп в схеме baseline vs adaptive}

{\def\LTcaptype{none} % do not increment counter
\begin{longtable}[]{|%
>{\RaggedRight\arraybackslash}p{(\linewidth - 9\tabcolsep) * \real{0.5000}}|
  >{\RaggedRight\arraybackslash}p{(\linewidth - 9\tabcolsep) * \real{0.5000}}|}
\hline
\begin{minipage}[b]{\linewidth}\RaggedRight
Группа
\end{minipage} & \begin{minipage}[b]{\linewidth}\RaggedRight
Описание
\end{minipage} \\
\hline
\endfirsthead
\hline\multicolumn{2}{|l|}{\small Продолжение таблицы 4}\\
\hline
\begin{minipage}[b]{\linewidth}\RaggedRight
Группа
\end{minipage} & \begin{minipage}[b]{\linewidth}\RaggedRight
Описание
\end{minipage} \\
\hline
\endhead
\hline
\endfoot
\hline
\endlastfoot
A, контрольная & Базовый неадаптивный сценарий без использования
\texttt{SER} в \texttt{policy\ layer}. \\
B, экспериментальная & Адаптивный сценарий с применением правил,
описанных во второй главе. \\
\end{longtable}
}

Единицей анализа является один завершенный телефонный контакт, то есть
одна попытка диалогового взаимодействия, доведенная до
регламентированного исхода: целевое действие, отказ, эскалация,
досрочное завершение или иное допустимое завершение разговора. Такое
определение единицы анализа позволяет одинаково учитывать как
результативные, так и конфликтные либо технически неуспешные сценарии
взаимодействия.

Для обеспечения внутренней валидности группы должны быть сопоставимы по
временным окнам запуска звонков, типу оффера, сегменту клиентской базы,
длине и структуре базового сценария, а также по используемому
технологическому стеку. Принципиально важно, чтобы единственным
существенным различием между группами выступал механизм эмоциональной
адаптации. Иначе зафиксированный эффект нельзя будет надежно
интерпретировать как следствие предложенной методики.

Независимой переменной в данном эксперименте выступает наличие либо
отсутствие emotion-aware \texttt{policy\ layer}. При значении \(0\)
сценарий не использует \texttt{SER} и не меняет тактику по
эмоциональному состоянию. При значении \(1\) сценарий использует
эмоциональную гипотезу, \texttt{confidence} и правила адаптации.
Зависимые переменные, в свою очередь, разделяются на конверсионные,
диалоговые, эмоционально-поведенческие и технические метрики.

Абляционные сравнения могут использоваться для последующей детализации
экспериментального дизайна. К числу таких сравнений относятся адаптация
по эмоции без учета \texttt{confidence}, адаптация по эмоции и
\texttt{confidence}, адаптация только по акустическому \texttt{SER} и
варианты частичной текстовой корректировки в неоднозначных случаях. В
настоящей работе основной эксперимент ограничивается сопоставлением
baseline- и adaptive-сценариев. Включение абляций методически оправдано,
поскольку современные работы по emotionally adaptive SDS также
используют сравнение с baseline и частичное отключение отдельных
компонентов политики {[}1; 5{]}. Общий дизайн эксперимента представлен
на рисунке 3.1.

\begin{figure}[H]
\centering
\VkrFigureExperimentDesign{\includegraphics[keepaspectratio,alt={Рисунок 3.1 --- Экспериментальный дизайн baseline vs adaptive с единым пулом контактов, общими исходами и системой метрик}]{docx_pipeline/figures/png/ch3_ab_experiment_design.png}}
\caption{Рисунок 3.1 --- Экспериментальный дизайн baseline vs adaptive с
единым пулом контактов, общими исходами и системой метрик}
\end{figure}
\FloatBarrier

\subsection{3.4 Система метрик и методика экспертной
оценки}\label{ux441ux438ux441ux442ux435ux43cux430-ux43cux435ux442ux440ux438ux43a-ux438-ux43cux435ux442ux43eux434ux438ux43aux430-ux44dux43aux441ux43fux435ux440ux442ux43dux43eux439-ux43eux446ux435ux43dux43aux438}

Для оценки вклада адаптивного модуля использовалась не одна итоговая
конверсионная величина, а согласованная система метрик. Это
принципиально важно для настоящей работы: рост целевых исходов не должен
достигаться ценой давления на абонента, роста конфликтности или
ухудшения качества диалога.

Основной конверсионной метрикой стала доля \texttt{target\_\allowbreak{}success}, то
есть доля разговоров, в которых абонент явно согласился на проверяемый
следующий шаг: получение коммерческого предложения или информации в
конкретный канал, согласованный повторный звонок, перевод на оператора
либо передачу контакта ответственному лицу. Простое предложение
менеджера отправить материалы без явного согласия клиента не считалось
успехом. Формулировки вида «подумаю», «может быть потом» или
«перезвоните как-нибудь» также не включались в \texttt{target\_\allowbreak{}success},
если следующий шаг не был проверяемым.

Дополнительно оценивались качество общения, доведение до следующего
шага, жесткие отказы, конфликтность и выраженность негативной реакции.
Такая сетка позволяет отделить прикладно полезный эффект адаптации от
формального роста конверсии, полученного за счет навязчивого или
небезопасного поведения системы.

\subsubsection{3.4.1 Схема кодирования экспертной
оценки}\label{ux441ux445ux435ux43cux430-ux43aux43eux434ux438ux440ux43eux432ux430ux43dux438ux44f-ux44dux43aux441ux43fux435ux440ux442ux43dux43eux439-ux43eux446ux435ux43dux43aux438}

Исходным фактическим материалом эксперимента выступали записи звонков,
полученные из телефонии компании. Для целей исследования эти записи были
транскрибированы и затем вручную оценены двумя экспертами по единой
схеме кодирования экспертных оценок. Итоговые числовые показатели,
использованные в расчетах, представлены как средние значения двух
экспертных оценок. Для каждого звонка фиксировались экспериментальная
группа, кампания, валидность контакта, итоговый исход, бинарный признак
\texttt{target\_\allowbreak{}success}, оценка доведения до следующего шага, пять
компонент качества общения, негативность, конфликтность и краткое
evidence-обоснование. Поля схемы экспертной оценки дополнительно
приведены в приложении А. Шкала оценки качества общения приведена в
таблице 5.

\VkrTableCaption{Таблица 5 --- Шкала экспертной оценки качества общения}

{\def\LTcaptype{none} % do not increment counter
\begin{longtable}[]{|%
>{\RaggedRight\arraybackslash}p{(\linewidth - 11\tabcolsep) * \real{0.2727}}|
  >{\centering\arraybackslash}p{(\linewidth - 11\tabcolsep) * \real{0.4545}}|
  >{\RaggedRight\arraybackslash}p{(\linewidth - 11\tabcolsep) * \real{0.2727}}|}
\hline
\begin{minipage}[b]{\linewidth}\RaggedRight
Компонента
\end{minipage} & \begin{minipage}[b]{\linewidth}\centering
Диапазон
\end{minipage} & \begin{minipage}[b]{\linewidth}\RaggedRight
Смысл оценки 2
\end{minipage} \\
\hline
\endfirsthead
\hline\multicolumn{3}{|l|}{\small Продолжение таблицы 5}\\
\hline
\begin{minipage}[b]{\linewidth}\RaggedRight
Компонента
\end{minipage} & \begin{minipage}[b]{\linewidth}\centering
Диапазон
\end{minipage} & \begin{minipage}[b]{\linewidth}\RaggedRight
Смысл оценки 2
\end{minipage} \\
\hline
\endhead
\hline
\endfoot
\hline
\endlastfoot
Релевантность & 0-2 & Ответ соответствует реплике клиента и стадии
диалога. \\
Ясность & 0-2 & Цель сформулирована кратко и без лишних повторов. \\
Эмпатия & 0-2 & Система признает состояние клиента и снижает
давление. \\
Работа с возражением & 0-2 & Возражение обработано корректно либо
разговор завершен спокойно. \\
Безопасность & 0-2 & Отсутствуют давление, манипуляции и продолжение
после жесткого отказа. \\
\end{longtable}
}

Итоговая интегральная оценка качества общения для звонка \(i\) после
усреднения экспертных оценок определялась как сумма пяти компонент:

\begin{equation}
q_i = r_i + c_i + e_i + o_i + b_i,
\quad q_i \in [0; 10],
\end{equation}

где \(r_i\) --- релевантность, \(c_i\) --- ясность, \(e_i\) --- эмпатия,
\(o_i\) --- обработка возражения, \(b_i\) --- безопасность общения.
Отдельно вводилась шкала доведения до следующего шага, представленная в
таблице 6.

\VkrTableCaption{Таблица 6 --- Шкала доведения разговора до следующего шага}

{\def\LTcaptype{none} % do not increment counter
\begin{longtable}[]{|c|l|}
\hline
Балл & Интерпретация \\
\hline
\endfirsthead
\hline\multicolumn{2}{|l|}{\small Продолжение таблицы 6}\\
\hline
Балл & Интерпретация \\
\hline
\endhead
\hline
\endfoot
\hline
\endlastfoot
0 & Разговор не дошел до релевантности или следующего шага. \\
1 & Частично выявлены релевантность, интерес или ответственное лицо. \\
2 & Следующий шаг предложен, но клиент не согласился. \\
3 & Клиент согласился в общем виде, но шаг не полностью проверяем. \\
4 & Следующий шаг явно согласован и проверяем. \\
\end{longtable}
}

Для присвоения поля \texttt{target\_\allowbreak{}success} равным \(1\) требовалось
значение поля \texttt{next\_\allowbreak{}step\_\allowbreak{}score}, равное \(4\). Таким образом,
целевой исход был связан не с субъективным впечатлением от разговора, а
с проверяемым действием клиента.

\subsubsection{3.4.2 Формальная постановка статистической
проверки}\label{ux444ux43eux440ux43cux430ux43bux44cux43dux430ux44f-ux43fux43eux441ux442ux430ux43dux43eux432ux43aux430-ux441ux442ux430ux442ux438ux441ux442ux438ux447ux435ux441ux43aux43eux439-ux43fux440ux43eux432ux435ux440ux43aux438}

Для основной конверсионной метрики вводилась бинарная случайная величина
\(Y_i\), принимающая значение \(1\) при достижении
\texttt{target\_\allowbreak{}success} и значение \(0\) в противном случае. В основном
анализе эта величина рассчитывалась не по всей исходной базе попыток
дозвона, а по аналитической выборке содержательных звонков длительностью
более \(40\) секунд. Для группы \(g \in \{A, B\}\) оценка доли целевых
исходов рассчитывалась как

\begin{equation}
\hat p_g^{(>40)} = \frac{s_g}{n_g},
\end{equation}

где \(n_g\) --- число звонков в выбранном аналитическом срезе, а \(s_g\)
--- число целевых исходов. Далее для краткости обозначение \(\hat p_g\)
используется именно для этой условной оценки по выборке \(T_i > 40\)
секунд.

Для связи с полной базой обзвона отдельно фиксировалась доля недозвонов:

\begin{equation}
r_{\mathrm{nd}} = \frac{N_{\mathrm{nd}}}{N_0} = 0{,}65,
\end{equation}

где \(N_0\) --- число попыток обзвона по базе, а \(N_{\mathrm{nd}}\) ---
число попыток, завершившихся недозвоном. Если требуется пересчитать
результат на всю базу попыток, необходимо учитывать не только недозвон,
но и долю состоявшихся контактов, попавших в аналитическую выборку
\(T_i > 40\) секунд:

\begin{equation}
p_g^{\mathrm{base}} =
(1-r_{\mathrm{nd},g}) \cdot \alpha_g \cdot \hat p_g^{(>40)},
\end{equation}

где \(\alpha_g\) --- доля состоявшихся контактов группы \(g\),
длительность которых превысила \(40\) секунд. В настоящей главе основное
сравнение выполняется по \(\hat p_g^{(>40)}\), поскольку именно в этой
выборке присутствует достаточный диалоговый материал для оценки качества
общения и доведения до целевого действия.

Основные меры эффекта задавались следующим образом:

\begin{equation}
\Delta = \hat p_B - \hat p_A,
\quad
L = \frac{\hat p_B - \hat p_A}{\hat p_A},
\quad
RR = \frac{\hat p_B}{\hat p_A}.
\end{equation}

Здесь \(\Delta\) показывает абсолютное изменение в долях, \(L\) ---
относительный прирост относительно контрольной группы, а \(RR\) ---
отношение вероятностей целевого исхода. Для среднего качества общения
использовалась величина

\begin{equation}
\bar q_g = \frac{1}{n_g}\sum_{i=1}^{n_g} q_i.
\end{equation}

Статистическая проверка основной гипотезы формулировалась как
одностороннее сравнение долей:

\begin{equation}
H_0: p_A = p_B,
\qquad
H_1: p_B > p_A.
\end{equation}

С учетом умеренного размера выборки и малых чисел успехов в отдельных
кампаниях применялся точный критерий Фишера. При фиксированном общем
числе успехов \(S = s_A + s_B\) число успехов в группе \(A\) имеет
гипергеометрическое распределение, а одностороннее \(p\)-значение для
альтернативы \(p_B > p_A\) рассчитывается как

\begin{equation}
p_{\mathrm{Fisher}} =
\sum_{x=0}^{s_A}
\frac{\binom{n_A}{x}\binom{n_B}{S-x}}
{\binom{n_A+n_B}{S}}.
\end{equation}

Для качества общения дополнительно использовался перестановочный тест по
статистике разности средних:

\begin{equation}
\Delta_q = \bar q_B - \bar q_A.
\end{equation}

Эта оценка является порядково-суммарной и не требует предположения о
нормальности распределения.

\subsection{3.5 Фактическая выборка и подготовка
данных}\label{ux444ux430ux43aux442ux438ux447ux435ux441ux43aux430ux44f-ux432ux44bux431ux43eux440ux43aux430-ux438-ux43fux43eux434ux433ux43eux442ux43eux432ux43aux430-ux434ux430ux43dux43dux44bux445}

При интерпретации эксперимента необходимо различать исходную базу
обзвона, состоявшиеся контакты и аналитическую выборку для оценки
качества диалога. По исходной базе, по которой выполнялся обзвон, около
\(65\%\) попыток завершались недозвоном. Эти случаи важны для общей
воронки кампании, но не содержат диалогового взаимодействия, поэтому не
позволяют оценивать качество общения, обработку возражений и доведение
клиента до следующего шага.

Для содержательного A/B-анализа использовалось отдельное правило отбора:
в аналитическую выборку включались все звонки, длительность которых
превышала \(40\) секунд и для которых был доступен транскрипт. Такой
порог выбран как прикладной критерий минимального содержательного
разговора: за меньший интервал обычно не успевают устойчиво проявиться
стадия предложения, реакция клиента, обработка возражения и фиксация
следующего шага.

Фактический массив аналитических данных включал 267 звонков
длительностью более \(40\) секунд, разбитых по трем кампаниям и двум
группам. Группа \texttt{A} соответствовала baseline-сценарию, а группа
\texttt{B} --- adaptive-сценарию с использованием адаптивного модуля.
Для каждого звонка были доступны аудиозапись и транскрипт, полученный
ASR-контуром. Состав набора экспериментальных звонков приведен в таблице
7.

Названия кампаний в таблицах и на рисунках приведены в виде обезличенных
служебных кодов \texttt{leadnet}, \texttt{merch} и \texttt{tprint}.
Такой способ обозначения используется для сохранения конфиденциальности
коммерческих офферов и партнерских данных; коды различают
экспериментальные сегменты, но не раскрывают реальные названия компаний,
клиентов или продуктов.

\VkrTableCaption{Таблица 7 --- Состав фактического набора экспериментальных звонков}

{\def\LTcaptype{none} % do not increment counter
\begin{longtable}[]{|l|r|r|r|}
\hline
Кампания & Группа A & Группа B & Всего \\
\hline
\endfirsthead
\hline\multicolumn{4}{|l|}{\small Продолжение таблицы 7}\\
\hline
Кампания & Группа A & Группа B & Всего \\
\hline
\endhead
\hline
\endfoot
\hline
\endlastfoot
leadnet & 46 & 46 & 92 \\
merch & 41 & 41 & 82 \\
tprint & 46 & 47 & 93 \\
Итого & 133 & 134 & 267 \\
\end{longtable}
}

Для статистического сравнения использовался основной срез
\texttt{primary}, из которого исключались несостоявшиеся контакты
\texttt{no\_\allowbreak{}contact} и технически непригодные случаи
\texttt{technical\_\allowbreak{}failure}. Это решение необходимо, поскольку цель
главы состоит в сравнении качества и результативности содержательного
диалога, а не в оценке качества дозвона или технической доступности
телефонной базы.

После очистки основной аналитический срез составил 227 разговоров: 112 в
группе \texttt{A} и 115 в группе \texttt{B}. При этом исходные 267
звонков сохранялись как полный raw-срез для контроля структуры данных и
возможного последующего аудита. Воронка отбора данных показана в таблице
8.

\VkrTableCaption{Таблица 8 --- Воронка отбора данных для экспериментального анализа}

{\def\LTcaptype{none} % do not increment counter
\begin{longtable}[]{|%
>{\RaggedRight\arraybackslash}p{(\linewidth - 11\tabcolsep) * \real{0.2727}}|
  >{\centering\arraybackslash}p{(\linewidth - 11\tabcolsep) * \real{0.4545}}|
  >{\RaggedRight\arraybackslash}p{(\linewidth - 11\tabcolsep) * \real{0.2727}}|}
\hline
\begin{minipage}[b]{\linewidth}\RaggedRight
Уровень
\end{minipage} & \begin{minipage}[b]{\linewidth}\centering
Значение
\end{minipage} & \begin{minipage}[b]{\linewidth}\RaggedRight
Назначение
\end{minipage} \\
\hline
\endfirsthead
\hline\multicolumn{3}{|l|}{\small Продолжение таблицы 8}\\
\hline
\begin{minipage}[b]{\linewidth}\RaggedRight
Уровень
\end{minipage} & \begin{minipage}[b]{\linewidth}\centering
Значение
\end{minipage} & \begin{minipage}[b]{\linewidth}\RaggedRight
Назначение
\end{minipage} \\
\hline
\endhead
\hline
\endfoot
\hline
\endlastfoot
Исходная база обзвона & 100\% попыток & Полный верхний уровень
кампании. \\
Недозвоны & 65\% попыток & Попытки без состоявшегося диалога. \\
Состоявшиеся контакты & около 35\% попыток & Потенциальная база для
диалогового анализа. \\
Raw \(T_i > 40\) секунд & 267 звонков & Все звонки с достаточной
длительностью и транскриптом. \\
Primary & 227 звонков & Содержательные разговоры без
\texttt{no\_\allowbreak{}contact} и \texttt{technical\_\allowbreak{}failure}. \\
\end{longtable}
}

Переход от raw-среза к primary-срезу по группам приведен в таблице 9.

\VkrTableCaption{Таблица 9 --- Переход от raw-среза к primary-срезу по группам}

{\def\LTcaptype{none} % do not increment counter
\begin{longtable}[]{|l|c|c|l|}
\hline
Срез & Группа A & Группа B & Назначение \\
\hline
\endfirsthead
\hline\multicolumn{4}{|l|}{\small Продолжение таблицы 9}\\
\hline
Срез & Группа A & Группа B & Назначение \\
\hline
\endhead
\hline
\endfoot
\hline
\endlastfoot
Raw \(T_i > 40\) секунд & 133 & 134 & Все аналитические звонки. \\
Primary & 112 & 115 & Разговоры без недозвонов и сбоев. \\
\end{longtable}
}

Итоговый исход каждого разговора выбирался из фиксированной таксономии:
\texttt{target\_\allowbreak{}success}, \texttt{qualified\_\allowbreak{}no\_\allowbreak{}commitment},
\texttt{soft\_\allowbreak{}refusal}, \texttt{hard\_\allowbreak{}refusal}, \texttt{no\_\allowbreak{}contact} и
\texttt{technical\_\allowbreak{}failure}. Такое разграничение позволило отдельно
учитывать целевое действие, частичный интерес без фиксации шага, мягкий
отказ, жесткий отказ и технически невалидные случаи.

\subsection{3.6 Результаты
A/B-сравнения}\label{ux440ux435ux437ux443ux43bux44cux442ux430ux442ux44b-ab-ux441ux440ux430ux432ux43dux435ux43dux438ux44f}

Основной результат сравнения заключается в том, что группа \texttt{B}
показала более высокую долю целевых исходов по primary-срезу: 29
успешных разговоров из 115 против 20 из 112 в группе \texttt{A}.

\begin{equation}
\hat p_A = \frac{20}{112} = 0{,}1786,
\qquad
\hat p_B = \frac{29}{115} = 0{,}2522.
\end{equation}

Абсолютный прирост составил

\begin{equation}
\Delta = 0{,}2522 - 0{,}1786 = 0{,}0736,
\end{equation}

то есть \(7{,}36\) процентного пункта. Относительный прирост по
отношению к контрольной группе составил

\begin{equation}
L = \frac{0{,}2522 - 0{,}1786}{0{,}1786} = 0{,}412,
\end{equation}

или около \(41{,}2\%\). Отношение долей целевого исхода равно

\begin{equation}
RR = \frac{0{,}2522}{0{,}1786} = 1{,}412.
\end{equation}

Иными словами, в пилотном эксперименте целевой исход в группе \texttt{B}
наблюдался примерно в \(1{,}41\) раза чаще, чем в группе \texttt{A}. При
этом точный критерий Фишера для односторонней альтернативы \(p_B > p_A\)
дал \(p = 0{,}1176\). Следовательно, результат следует интерпретировать
как положительный пилотный эффект и практически значимый сигнал, но не
как статистически доказанное превосходство на уровне значимости
\(0{,}05\). Primary-сводка по группам и кампаниям приведена в таблице
10.

\VkrTableCaption{Таблица 10 --- Primary-сводка по группам и кампаниям}

{\def\LTcaptype{none} % do not increment counter
\begin{longtable}[]{|l|l|r|r|r|r|r|}
\hline
Кампания & Группа & N & Успехи & Доля успеха & Качество & Next step \\
\hline
\endfirsthead
\hline\multicolumn{7}{|l|}{\small Продолжение таблицы 10}\\
\hline
Кампания & Группа & N & Успехи & Доля успеха & Качество & Next step \\
\hline
\endhead
\hline
\endfoot
\hline
\endlastfoot
all & A & 112 & 20 & 17,9\% & 5,04 & 1,68 \\
all & B & 115 & 29 & 25,2\% & 5,21 & 1,83 \\
leadnet & A & 37 & 3 & 8,1\% & 4,57 & 0,92 \\
leadnet & B & 40 & 7 & 17,5\% & 4,08 & 1,35 \\
merch & A & 35 & 12 & 34,3\% & 5,80 & 2,37 \\
merch & B & 36 & 12 & 33,3\% & 6,22 & 2,08 \\
tprint & A & 40 & 5 & 12,5\% & 4,83 & 1,77 \\
tprint & B & 39 & 10 & 25,6\% & 5,44 & 2,10 \\
\end{longtable}
}

Для наглядного сопоставления долей целевого исхода по группам и
кампаниям результаты представлены на рисунке 3.2. Визуализация
показывает, что общий прирост adaptive-группы связан не с равномерным
улучшением по всем кампаниям, а с выраженной положительной динамикой в
кампаниях \texttt{leadnet} и \texttt{tprint} при почти неизменной доле
успеха в кампании \texttt{merch}.

\begin{figure}[H]
\centering
\VkrFigureInline{\includegraphics[keepaspectratio,alt={Рисунок 3.2 --- Доля целевых исходов target\_success в primary-срезе по группам и кампаниям}]{docx_pipeline/figures/png/ch3_target_success_rates.png}}
\caption{Рисунок 3.2 --- Доля целевых исходов target\_success в
primary-срезе по группам и кампаниям}
\end{figure}

Таблица 11 показывает, что общий рост в группе \texttt{B} сформирован
прежде всего за счет кампаний \texttt{leadnet} и \texttt{tprint}. В
кампании \texttt{merch} целевая доля осталась практически на том же
уровне, однако среднее качество общения в группе \texttt{B} оказалось
выше.

\VkrTableCaption{Таблица 11 --- Размер эффекта и статистическая проверка по primary-срезу}

{\def\LTcaptype{none} % do not increment counter
\begin{longtable}[]{|l|c|c|c|c|c|}
\hline
Кампания & N A/B & Δ, п.п. & RR & p Fisher & Δ качества \\
\hline
\endfirsthead
\hline\multicolumn{6}{|l|}{\small Продолжение таблицы 11}\\
\hline
Кампания & N A/B & Δ, п.п. & RR & p Fisher & Δ качества \\
\hline
\endhead
\hline
\endfoot
\hline
\endlastfoot
all & 112/115 & 7,36 & 1,41 & 0,1176 & 0,16 \\
leadnet & 37/40 & 9,39 & 2,16 & 0,1887 & -0,49 \\
merch & 35/36 & -0,95 & 0,97 & 0,6314 & 0,42 \\
tprint & 40/39 & 13,14 & 2,05 & 0,1145 & 0,61 \\
\end{longtable}
}

На рисунке 3.3 дополнительно показан размер эффекта как разность долей
\texttt{target\_\allowbreak{}success} между группой \texttt{B} и группой \texttt{A}.
Такой вид представления удобен для защиты, поскольку сразу показывает
направление эффекта и неоднородность результата по кампаниям.

\begin{figure}[H]
\centering
\VkrFigureCompact{\includegraphics[keepaspectratio,alt={Рисунок 3.3 --- Размер эффекта adaptive-сценария по кампаниям в primary-срезе}]{docx_pipeline/figures/png/ch3_effect_delta_by_campaign.png}}
\caption{Рисунок 3.3 --- Размер эффекта adaptive-сценария по кампаниям в
primary-срезе}
\end{figure}

По качеству общения общий прирост составил \(0{,}16\) балла по шкале
\(0\)-\(10\) (\(5{,}21\) против \(5{,}04\)). Перестановочный тест по
разности средних дал \(p = 0{,}2837\), поэтому этот рост также
рассматривается как направленный, но не статистически устойчивый на
уровне \(0{,}05\). Более выраженная положительная динамика качества
наблюдалась в кампаниях \texttt{merch} и \texttt{tprint}, тогда как в
\texttt{leadnet} среднее качество в группе \texttt{B} оказалось ниже,
несмотря на рост доли целевого исхода.

Распределение исходов подтверждает, что рост \texttt{target\_\allowbreak{}success} в
группе \texttt{B} сопровождался сокращением числа разговоров с частичным
интересом без фиксации следующего шага. Это означает, что часть диалогов
в adaptive-группе чаще доводилась до конкретного проверяемого действия,
а не оставалась в промежуточном состоянии. Распределение исходов
приведено в таблице 12.

\VkrTableCaption{Таблица 12 --- Распределение исходов в primary-срезе}

{\def\LTcaptype{none} % do not increment counter
\begin{longtable}[]{|l|r|r|}
\hline
Исход & Группа A & Группа B \\
\hline
\endfirsthead
\hline\multicolumn{3}{|l|}{\small Продолжение таблицы 12}\\
\hline
Исход & Группа A & Группа B \\
\hline
\endhead
\hline
\endfoot
\hline
\endlastfoot
\texttt{target\_\allowbreak{}success} & 20 & 29 \\
\texttt{qualified\_\allowbreak{}no\_\allowbreak{}commitment} & 22 & 14 \\
\texttt{soft\_\allowbreak{}refusal} & 48 & 49 \\
\texttt{hard\_\allowbreak{}refusal} & 22 & 23 \\
\end{longtable}
}

Та же структура исходов в долевом виде приведена на рисунке 3.4. Он
показывает, что в группе \texttt{B} увеличилась доля целевых исходов и
одновременно уменьшилась доля разговоров с частичным интересом без
фиксации следующего шага. При этом доли мягких и жестких отказов
остались близкими.

\begin{figure}[H]
\centering
\VkrFigureWideCompact{\includegraphics[keepaspectratio,alt={Рисунок 3.4 --- Распределение исходов разговора в primary-срезе по группам A и B}]{docx_pipeline/figures/png/ch3_outcome_distribution.png}}
\caption{Рисунок 3.4 --- Распределение исходов разговора в primary-срезе
по группам A и B}
\end{figure}

При этом число жестких отказов и конфликтных завершений в целом осталось
сопоставимым: 22 случая в группе \texttt{A} и 23 случая в группе
\texttt{B}. Следовательно, по имеющимся данным рост целевых исходов не
сопровождался явным увеличением конфликтности на общем уровне. В
кампании \texttt{merch} количество жестких отказов в группе \texttt{B}
было ниже, чем в \texttt{A} (\(6\) против \(8\)), тогда как в
\texttt{tprint} оно было выше (\(10\) против \(7\)), что указывает на
неоднородность эффекта по кампаниям.

\subsection{3.7 Интерпретация
результатов}\label{ux438ux43dux442ux435ux440ux43fux440ux435ux442ux430ux446ux438ux44f-ux440ux435ux437ux443ux43bux44cux442ux430ux442ux43eux432}

Полученные данные в целом согласуются с основной гипотезой исследования:
adaptive-сценарий показал более высокую долю доведения разговора до
целевого действия. Абсолютный прирост \(7{,}36\) процентного пункта и
относительный прирост около \(41{,}2\%\) являются практически значимыми
для прикладного автообзвона, поскольку даже умеренное увеличение доли
проверяемых следующих шагов может существенно влиять на экономику
кампании.

Вместе с тем статистическая интерпретация должна оставаться осторожной.
Односторонний точный критерий Фишера дал \(p = 0{,}1176\), то есть
наблюдаемое различие не достигает традиционного уровня значимости
\(0{,}05\). Поэтому результат корректно формулировать не как
окончательно доказанное превосходство адаптивного модуля, а как
положительный пилотный сигнал, подтверждающий целесообразность
дальнейшего масштабирования эксперимента. Сила этого результата состоит
не в формальном статистическом доказательстве, а в величине наблюдаемого
эффекта и его прикладной интерпретируемости: прирост более чем на семь
процентных пунктов заметен для бизнес-воронки, но требует подтверждения
на большей выборке и при сохранении той же процедуры выгрузки данных из
телефонии и экспертной оценки транскриптов.

По отдельным кампаниям эффект оказался неоднородным. В \texttt{leadnet}
доля целевых исходов в группе \texttt{B} выросла с \(8{,}1\%\) до
\(17{,}5\%\), но средняя экспертная оценка качества общения снизилась.
Это может означать, что adaptive-сценарий чаще фиксировал следующий шаг,
но не всегда делал это с более высоким качеством диалога. В
\texttt{merch} конверсионная доля практически не изменилась, однако
качество общения выросло, а число жестких отказов снизилось. В
\texttt{tprint} наблюдалось одновременно увеличение целевых исходов и
рост среднего качества, что делает эту кампанию наиболее благоприятным
примером работы адаптивного модуля.

Содержательно это означает, что эффект эмоциональной адаптации не
является универсальной константой. Он зависит от типа оффера, структуры
клиентской базы, содержания возражений и того, насколько часто в
разговоре возникает ситуация, где изменение тактики ответа действительно
влияет на исход. Поэтому предложенная методика должна рассматриваться
как управляемый контур адаптации, требующий последующей калибровки под
конкретный сценарий, а не как универсальный способ механически увеличить
конверсию во всех кампаниях.

\subsection{3.8 Угрозы валидности и
ограничения}\label{ux443ux433ux440ux43eux437ux44b-ux432ux430ux43bux438ux434ux43dux43eux441ux442ux438-ux438-ux43eux433ux440ux430ux43dux438ux447ux435ux43dux438ux44f}

Основной источник данных эксперимента составили записи звонков,
полученные из телефонии компании. Транскрибирование, группировка
материалов по кампаниям и экспериментальным условиям, исключение
непригодных контактов и итоговая экспертная разметка выполнялись уже в
рамках настоящего исследования. Такой контур достаточен для задачи ВКР,
поскольку проверяется различие между двумя сценариями автообзвона в
пределах одного и того же телефонного процесса. Поэтому
\texttt{target\_\allowbreak{}success} в настоящем анализе трактуется как экспертно
выявленное согласие на следующий проверяемый шаг внутри разговора, а не
как отдельная коммерческая метрика, требующая раскрытия конфиденциальных
внутренних данных партнера.

Второе ограничение связано с формированием аналитической выборки. По
исходной базе обзвона около \(65\%\) попыток завершались недозвоном, а
содержательный анализ проводился только по звонкам длительностью более
\(40\) секунд. Поэтому полученные оценки характеризуют не всю базу
попыток дозвона, а условный эффект adaptive-сценария среди разговоров,
где состоялось достаточно длительное диалоговое взаимодействие. Для
расчета сквозной конверсии по всей базе необходим полный лог дозвона с
числом попыток, недозвонов, коротких контактов и звонков длительностью
более \(40\) секунд по каждой группе.

Третье ограничение связано с размером выборки. Хотя аналитический массив
включал 267 звонков длительностью более \(40\) секунд, основной
primary-срез содержал 227 содержательных разговоров, а подвыборки по
отдельным кампаниям были еще меньше. Именно поэтому для основной доли
использовался точный критерий Фишера, а полученный результат трактуется
как пилотный.

Четвертое ограничение связано с источником разметки. Ручная экспертная
оценка повышает содержательную точность интерпретации звонков, но
остается зависимой от схемы кодирования оценок и согласованности
оценщиков. В данной пилотной проверке использовались оценки двух
экспертов, а итоговые числовые показатели рассчитывались как среднее
между ними. Это снижает зависимость результата от одного оценщика,
однако не заменяет формальной проверки межэкспертного согласия,
поскольку коэффициенты согласованности отдельно не рассчитывались.
Поэтому результаты ручной оценки следует интерпретировать как
исследовательскую экспертную разметку, достаточную для первичной
проверки гипотезы, но требующую дополнительной валидации при
масштабировании эксперимента. В последующих циклах целесообразно
рассчитывать каппу Коэна для бинарных и категориальных признаков, таких
как \texttt{target\_\allowbreak{}success}, \texttt{outcome} и \texttt{conflict}, а
также взвешенную каппу или \texttt{ICC} для порядковых шкал
\texttt{next\_\allowbreak{}step\_\allowbreak{}score}, \texttt{communication\_\allowbreak{}quality\_\allowbreak{}score} и
\texttt{negativity\_\allowbreak{}score}.

Пятое ограничение состоит в неполноте технических логов. В доступном
массиве основным наблюдаемым объектом были аудиозапись и транскрипт;
поэтому часть технических метрик, связанных с \texttt{SER\ confidence},
задержками пайплайна, \texttt{fallback}-переходами и внутренними
решениями \texttt{policy\ layer}, не могла быть рассчитана с той же
степенью надежности, что экспертные метрики качества диалога. Эти
показатели остаются обязательными для масштабирования экспериментального
контура, где логирование должно вестись на уровне отдельного такта.

Наконец, внешняя валидность ограничена тремя конкретными кампаниями:
\texttt{leadnet}, \texttt{merch} и \texttt{tprint}. Перенос результата
на другие офферы, отрасли и сегменты клиентской базы требует отдельной
проверки. Это особенно важно с учетом работ по русскоязычному SER, где
качество моделей чувствительно к языку, корпусу, каналу записи и
структуре речевого материала {[}6; 21{]}.

\subsection{3.9 Выводы по
главе}\label{ux432ux44bux432ux43eux434ux44b-ux43fux43e-ux433ux43bux430ux432ux435-2}

В данной главе сформирован и применен воспроизводимый контур
экспериментальной проверки предложенной методики адаптивного
автообзвона. Зафиксированы гипотезы, экспериментальный дизайн
\texttt{baseline\ vs\ adaptive}, схема ручной экспертной оценки,
формальные метрики, статистическая проверка и ограничения интерпретации.

Фактический pilot A/B-анализ 267 звонков показал, что после исключения
несостоявшихся и технически непригодных контактов adaptive-группа
\texttt{B} достигла доли \texttt{target\_\allowbreak{}success} \(25{,}2\%\) против
\(17{,}9\%\) у baseline-группы \texttt{A}. Абсолютный прирост составил
\(7{,}36\) процентного пункта, относительный прирост --- около
\(41{,}2\%\), а отношение долей успеха --- \(1{,}41\).

Одновременно среднее качество общения в группе \texttt{B} оказалось
немного выше (\(5{,}21\) против \(5{,}04\) по шкале \(0\)--\(10\)), а
общее число жестких отказов и конфликтных завершений осталось
сопоставимым. Статистическая проверка не позволяет утверждать значимость
эффекта на уровне \(0{,}05\), однако результаты дают положительный
пилотный сигнал в пользу адаптивного модуля и подтверждают практическую
целесообразность дальнейшего масштабирования эксперимента с более полным
техническим логированием тактов диалога.

\VkrStructuralSection{ЗАКЛЮЧЕНИЕ}\label{ux437ux430ux43aux43bux44eux447ux435ux43dux438ux435}

Основные результаты исследования состоят в следующем. В настоящей работе
решена задача формализации методики адаптивного автообзвона,
ориентированной на русскоязычный исходящий телефонный диалог и
учитывающей эмоциональное состояние собеседника как управляемый параметр
выбора следующего действия системы. В отличие от подходов, в которых
эмоция влияет главным образом на стиль ответа, предложенная методика
задает явную связь между состоянием абонента, надежностью его
распознавания, стадией сценария и допустимой тактикой разговора.

В ходе работы сформирован воспроизводимый архитектурный контур,
включающий \texttt{VAD}, \texttt{ASR}, \texttt{SER}, \texttt{LLM},
\texttt{TTS}, слой нормализации и rule-based \texttt{policy\ layer}. Для
каждого такта диалога зафиксированы входные данные, выходные сигналы и
правила перехода к следующему действию системы. Тем самым авторский
результат работы заключается не только в выборе конкретного набора
моделей, но прежде всего в формальном описании их взаимодействия в
составе единой методики.

Содержательным ядром предложенного решения выступает ограниченный
рабочий набор эмоциональных состояний, включающий нейтральность,
заинтересованность, раздражение, злость и неопределенность. Для этих
состояний сформирована таблица policy-правил, связывающая эмоциональный
режим собеседника с темпом ответа, степенью эмпатии, глубиной
аргументации, предложением альтернативы и условиями прекращения
автоматического убеждения. Дополнительно показано, что параметр
\texttt{confidence} должен рассматриваться не как вспомогательная
статистика, а как обязательный управляющий параметр, ограничивающий
глубину адаптации и снижающий риск ложных критических переходов.

Отдельно следует выделить научный и практический результаты работы.
Научный результат настоящего этапа состоит в переходе от общего тезиса о
полезности emotion-aware систем к формализованной и воспроизводимой
методике, пригодной для прикладного русскоязычного автообзвона. По
результатам анализа литературы показано, что открытые исследования
подтверждают перспективность учета эмоций в разговорных системах, однако
не дают завершенного методического контура, одновременно объединяющего
\texttt{SER}, оценку надежности, rule-based \texttt{policy\ layer},
\texttt{LLM} как языкового исполнителя и controlled
\texttt{baseline\ vs\ adaptive} эксперимент. Тем самым работа закрывает
выявленный исследовательский разрыв между современными emotion-aware
подходами и практической задачей регламентированного исходящего обзвона.

Практический результат состоит в том, что методика изначально
проектируется как объяснимый инженерный контур, пригодный для
развертывания в инфраструктуре агентства или контактного центра. Для
этой цели зафиксированы требования к журналированию тактов диалога,
правила \texttt{fallback}-поведения, условия эскалации на оператора, а
также логика разделения baseline- и adaptive-сценариев. Такой подход
позволяет анализировать не только итоговую конверсию, но и конкретные
причины переходов, отказов, конфликтов и безопасных завершений
разговора.

Существенным практическим шагом исследования стала фиксация пилотного
сценария исходящего звонка как структурного use case, в котором цель
системы состоит в согласовании следующего содержательного контакта. На
этой основе была проведена ручная экспертная оценка 267
экспериментальных звонков, разделенных на baseline-группу \texttt{A} и
adaptive-группу \texttt{B}. По primary-срезу без несостоявшихся и
технически непригодных контактов группа \texttt{B} показала более
высокую долю \texttt{target\_\allowbreak{}success}: \(25{,}2\%\) против \(17{,}9\%\)
у группы \texttt{A}. Абсолютный прирост составил \(7{,}36\) процентного
пункта, относительный прирост --- около \(41{,}2\%\), а отношение долей
успеха --- \(1{,}41\).

Полученный результат следует трактовать как положительный пилотный
сигнал, а не как окончательно доказанное статистическое превосходство:
односторонний точный критерий Фишера по основной доле дал
\(p = 0{,}1176\). Тем не менее эксперимент показал, что предложенная
методика может улучшать доведение разговора до проверяемого следующего
шага без явного роста общей конфликтности.

При интерпретации полученных результатов необходимо учитывать
ограничения исследования. Они связаны прежде всего с пилотным характером
экспериментальной проверки. По исходной базе обзвона около \(65\%\)
попыток завершались недозвоном, а содержательная экспертная оценка
выполнялась по выборке звонков длительностью более \(40\) секунд.
Поэтому полученный эффект характеризует качество и результативность
adaptive-сценария внутри достаточно длинных состоявшихся разговоров, но
не является полной сквозной конверсией по всей базе попыток дозвона.

Фактический массив достаточен для первичной оценки направления эффекта,
но не позволяет утверждать статистически значимое превосходство
adaptive-модуля на уровне \(0{,}05\). При этом \texttt{target\_\allowbreak{}success}
оценивался по содержанию телефонного разговора и подготовленного в
рамках исследования транскрипта как согласие на следующий проверяемый
шаг. Такой уровень разметки соответствует цели ВКР: сравнить baseline- и
adaptive-сценарии внутри одного телефонного процесса без раскрытия
конфиденциальных коммерческих данных партнера.

Кроме того, перенос предложенной методики на иные офферы и иные сегменты
клиентской базы требует отдельной проверки. Хотя общий принцип emotion +
confidence + stage → action сохраняет универсальность, фактический
эффект адаптации может зависеть от отрасли, характера предложения,
качества аудиоканала, структуры клиентской базы и организационных
ограничений контактного центра. Поэтому внешняя валидность результатов
должна оцениваться с учетом конкретного пилотного сценария.

Необходимо учитывать и то обстоятельство, что качество итогового
поведения системы определяется не одним модулем \texttt{SER}, а всей
каскадной цепочкой \texttt{VAD\ +\ ASR\ +\ SER\ +\ LLM\ +\ TTS}. Ошибки
распознавания речи, нестабильность эмоциональной классификации, задержки
пайплайна или непрохождение ответа через валидацию способны снижать
прикладной эффект даже при корректно сформулированной методике. Именно
поэтому в работе отдельно выделены требования к логированию, контролю
латентности и анализу \texttt{fallback}-переходов.

Дальнейшее развитие работы связано с масштабированием A/B-эксперимента и
расширением технического логирования до уровня каждого такта диалога.
При масштабировании необходимо фиксировать не только итоговый транскрипт
и outcome, но и \texttt{SER\ confidence}, нормализованное эмоциональное
состояние, выбранное \texttt{policy}-действие,
\texttt{fallback}-переходы, задержки модулей и факты эскалации на
оператора. Это позволит проверить не только итоговый эффект
adaptive-модуля, но и внутренний механизм, за счет которого достигается
улучшение.

Отдельным направлением развития является калибровка порогов
\(\theta_{\mathrm{low}}\) и \(\theta_{\mathrm{high}}\) на расширенной
выборке, а также сравнение вариантов policy layer: без учета эмоций, с
учетом только класса эмоции, с учетом эмоции и \texttt{confidence}, а
также с дополнительной текстовой валидацией спорных случаев.

Итоговый вывод состоит в том, что в работе сформирована целостная
исследовательская, методическая и экспериментальная основа ВКР: выполнен
аналитический обзор, разработана формальная методика адаптивного
автообзвона, описан baseline/adaptive-дизайн, проведена ручная
экспертная оценка фактических A/B-транскриптов и рассчитаны основные
метрики эффекта. Эксперимент не дает оснований для категорического
утверждения о статистически доказанном превосходстве adaptive-сценария,
но показывает практически значимый положительный пилотный эффект и
подтверждает обоснованность дальнейшего развития предложенной методики.

\VkrStructuralSection{СПИСОК ИСПОЛЬЗОВАННЫХ ИСТОЧНИКОВ}\label{ux441ux43fux438ux441ux43eux43a-ux438ux441ux43fux43eux43bux44cux437ux43eux432ux430ux43dux43dux44bux445-ux438ux441ux442ux43eux447ux43dux438ux43aux43eux432}

\begin{enumerate}
\def\labelenumi{\arabic{enumi}.}
\tightlist
\item
  Feng S. et al.~Infusing Emotions into Task-oriented Dialogue Systems
  // Proceedings of SIGDIAL 2024. 2024. URL:
  \url{https://aclanthology.org/2024.sigdial-1.60/} (дата обращения:
  12.04.2026).
\item
  Zhu Y. et al.~An emotion-sensitive dialogue policy for task-oriented
  dialogue system // Scientific Reports. 2024. URL:
  \url{https://www.nature.com/articles/s41598-024-70463-x} (дата
  обращения: 12.04.2026).
\item
  Zhang Y. et al.~Affective Computing in the Era of Large Language
  Models: A Survey {[}Электронный ресурс{]}. 2024. URL:
  \url{https://arxiv.org/abs/2408.04638} (дата обращения: 12.04.2026).
\item
  Xue S. et al.~E-chat: Emotion-sensitive Spoken Dialogue System with
  Large Language Models {[}Электронный ресурс{]}. 2024. URL:
  \url{https://arxiv.org/abs/2401.00475} (дата обращения: 12.04.2026).
\item
  Yoon J. et al.~EmoSDS: Unified Emotionally Adaptive Spoken Dialogue
  System // Future Internet. 2025. URL:
  \url{https://www.mdpi.com/1999-5903/17/4/143} (дата обращения:
  12.04.2026).
\item
  Лемаев В. И., Лукашевич Н. В. Автоматическая классификация эмоций в
  речи: методы и данные // Litera. 2024. URL:
  \url{https://journals.rcsi.science/2409-8698/article/view/379468}
  (дата обращения: 12.04.2026).
\item
  Martin-Donas J. M. et al.~Speech Emotion Recognition for Call Centers
  using Self-supervised Models // Proceedings of ICNLSP 2024. 2024. URL:
  \url{https://aclanthology.org/2024.icnlsp-1.14/} (дата обращения:
  12.04.2026).
\item
  Brun L. et al.~Exploring Emotion-Sensitive LLM-Based Conversational AI
  {[}Электронный ресурс{]}. 2025. URL:
  \url{https://arxiv.org/abs/2502.08920} (дата обращения: 12.04.2026).
\item
  Cui C. et al.~Recent Advances in Speech Language Models: A Survey
  {[}Электронный ресурс{]}. 2024. URL:
  \url{https://arxiv.org/abs/2410.03751} (дата обращения: 12.04.2026).
\item
  Rakib M. R. et al.~DialogXpert: Driving Intelligent and Emotion-Aware
  Conversations through Online Value-Based Reinforcement Learning with
  LLM Priors {[}Электронный ресурс{]}. 2025. URL:
  \url{https://arxiv.org/abs/2505.17795} (дата обращения: 12.04.2026).
\item
  Wang Z. et al.~Empathy Omni: Enabling Empathetic Speech Response
  Generation through Large Language Models {[}Электронный ресурс{]}.
  2025. URL: \url{https://arxiv.org/abs/2508.18655} (дата обращения:
  12.04.2026).
\item
  Cornell S. et al.~Towards Using TTS-Generated Context for Improving
  LLM-Based Conversational ASR {[}Электронный ресурс{]}. 2024. URL:
  \url{https://arxiv.org/abs/2408.09215} (дата обращения: 12.04.2026).
\item
  Veluri B. et al.~Beyond Turn-Based Interfaces: Synchronous LLMs as
  Full-Duplex Dialogue Agents // Proceedings of EMNLP 2024. 2024. URL:
  \url{https://aclanthology.org/2024.emnlp-main.1192/} (дата обращения:
  12.04.2026).
\item
  Wang C. et al.~Full-Duplex Speech Dialogue Scheme Based on Large
  Language Models {[}Электронный ресурс{]}. 2024. URL:
  \url{https://arxiv.org/abs/2405.19487} (дата обращения: 12.04.2026).
\item
  Jafarzadeh M. et al.~Optimizing Speech Emotion Recognition with
  Wav2Vec2 and HuBERT {[}Электронный ресурс{]}. 2024. URL:
  \url{https://arxiv.org/abs/2411.02964} (дата обращения: 12.04.2026).
\item
  Ma Z. et al.~emotion2vec: Self-Supervised Pre-Training for Speech
  Emotion Representation // Findings of ACL 2024. 2024. URL:
  \url{https://aclanthology.org/2024.findings-acl.931/} (дата обращения:
  12.04.2026).
\item
  Liu R. et al.~Emotion-Aware Self-Supervised Learning for Speech
  Emotion Recognition // Proceedings of Interspeech 2024. 2024. URL:
  \url{https://www.isca-archive.org/interspeech_2024/liu24r_interspeech.html}
  (дата обращения: 12.04.2026).
\item
  Li Y. et al.~Revise, Reason, and Recognize: LLM-Based Emotion
  Recognition from Human Speech with ASR Errors {[}Электронный
  ресурс{]}. 2024. URL: \url{https://arxiv.org/abs/2409.15551} (дата
  обращения: 12.04.2026).
\item
  Zhang Y. et al.~MERSA: A Large-Scale Multimodal Emotion Recognition in
  Speech Dataset // Proceedings of ACL 2024. 2024. URL:
  \url{https://aclanthology.org/2024.acl-long.752/} (дата обращения:
  12.04.2026).
\item
  Singh A. et al.~Hybrid Emotion Recognition for Customer Interactions:
  A Privacy-Aware Multimodal Approach {[}Электронный ресурс{]}. 2025.
  URL: \url{https://arxiv.org/abs/2503.21927} (дата обращения:
  12.04.2026).
\item
  Balabanova T. N., Gaivoronskaya D. I., Doborovich A. N. Распознавание
  эмоций по устной речи с использованием нейросетевого подхода //
  Research Result. Theoretical and Applied Linguistics. 2024. URL:
  \url{https://rrlinguistics.ru/en/journal/annotation/3576/} (дата
  обращения: 12.04.2026).
\item
  Mallol-Ragolta A., Schuller B. W. Coupling Sentiment and Arousal
  Analysis Towards an Affective Dialogue Manager // IEEE Access. 2024.
  URL: \url{https://doi.org/10.1109/ACCESS.2024.3361750} (дата
  обращения: 12.04.2026).
\end{enumerate}

\VkrAppendixSection{А}{Схема кодирования экспертной оценки экспериментальных звонков}\label{ux43fux440ux438ux43bux43eux436ux435ux43dux438ux435-ux430-ux441ux445ux435ux43cux430-ux43aux43eux434ux438ux440ux43eux432ux430ux43dux438ux44f-ux44dux43aux441ux43fux435ux440ux442ux43dux43eux439-ux43eux446ux435ux43dux43aux438-ux44dux43aux441ux43fux435ux440ux438ux43cux435ux43dux442ux430ux43bux44cux43dux44bux445-ux437ux432ux43eux43dux43aux43eux432}

Для ручной оценки A/B-транскриптов использовалась единая схема
кодирования экспертных оценок. Оценку выполняли два эксперта, а итоговые
числовые показатели представлялись как средние значения двух оценок.
Схема фиксировала валидность контакта, итоговый исход, достижение
целевого действия, качество общения, доведение до следующего шага,
негативность и конфликтность. Единицей анализа являлся один транскрипт
звонка. Основные поля схемы приведены в таблице А.1.

\VkrTableCaption{Таблица А.1 --- Основные поля схемы кодирования экспертной оценки}

{\def\LTcaptype{none} % do not increment counter
\begin{longtable}[]{|%
>{\RaggedRight\arraybackslash}p{(\linewidth - 11\tabcolsep) * \real{0.2727}}|
  >{\centering\arraybackslash}p{(\linewidth - 11\tabcolsep) * \real{0.4545}}|
  >{\RaggedRight\arraybackslash}p{(\linewidth - 11\tabcolsep) * \real{0.2727}}|}
\hline
\begin{minipage}[b]{\linewidth}\RaggedRight
Поле
\end{minipage} & \begin{minipage}[b]{\linewidth}\centering
Диапазон
\end{minipage} & \begin{minipage}[b]{\linewidth}\RaggedRight
Назначение
\end{minipage} \\
\hline
\endfirsthead
\hline\multicolumn{3}{|l|}{\small Продолжение таблицы А.1}\\
\hline
\begin{minipage}[b]{\linewidth}\RaggedRight
Поле
\end{minipage} & \begin{minipage}[b]{\linewidth}\centering
Диапазон
\end{minipage} & \begin{minipage}[b]{\linewidth}\RaggedRight
Назначение
\end{minipage} \\
\hline
\endhead
\hline
\endfoot
\hline
\endlastfoot
\texttt{valid\_\allowbreak{}contact} & 0/1 & Был ли содержательный обмен
репликами. \\
\texttt{outcome} & категория & Итог разговора по фиксированной
таксономии. \\
\texttt{target\_\allowbreak{}success} & 0/1 & Явное согласие на проверяемый следующий
шаг. \\
\texttt{next\_\allowbreak{}step\_\allowbreak{}score} & 0-4 & Степень доведения до конкретного
следующего действия. \\
\texttt{communication\_\allowbreak{}quality\_\allowbreak{}score} & 0-10 & Сумма пяти компонент
качества общения. \\
\texttt{negativity\_\allowbreak{}score} & 0-3 & Выраженность негативной реакции
клиента. \\
\texttt{conflict} & 0/1 & Наличие жесткого отказа или конфликтного
завершения. \\
\texttt{evidence} & текст & Короткое основание оценки по транскрипту. \\
\end{longtable}
}

Качество общения рассчитывалось как сумма пяти компонент по шкале
\(0\)--\(10\). Целевой исход фиксировался только при
\texttt{next\_step\_score\ =\ 4}, то есть при явно согласованном и
проверяемом следующем шаге.


\end{document}
